% Lectura y comentario : relaciones internacionales, relaciones Norte-Sur, sistemas políticos, derechos y deberes — Espagnol Tle A — Cours signé OnBuch+
% Programme officiel MINESEC. Compiler avec : tectonic E13-relaciones-internacionales-norte-sur.tex
\newcommand{\DOCMATIERE}{Espagnol}
\newcommand{\DOCNIVEAU}{Tle A}
\newcommand{\DOCMODULE}{Módulo 3 : Vie citoyenne et ouverture au monde}
\newcommand{\DOCLECON}{E13}
\newcommand{\DOCTITRE}{Lectura y comentario : relaciones internacionales, relaciones Norte-Sur, sistemas políticos, derechos y deberes}
% OnBuch+ — préambule « cours signés » ENRICHI (compilé avec tectonic / XeLaTeX)
% Système visuel « Pop » d'OnBuch+ : fond crème (jamais blanc), bordures encre
% épaisses, ombres « dures » décalées, titres Archivo Black en capitales.
% Version MATHÉMATIQUES : graphes (pgfplots/tikz), boîtes pédagogiques
% (définition, propriété, méthode, attention, le savais-tu, exercice résolu,
% exercice, TP, bilan), page de garde et signature OnBuch+ ancrée en bas.
%
% Macros attendues dans main.tex AVANT \input{preamble} :
%   \DOCMATIERE  \DOCNIVEAU  \DOCMODULE  \DOCLECON (numéro)  \DOCTITRE
\documentclass[11pt]{article}

\usepackage{fontspec}
\usepackage[spanish,french]{babel} % français = langue principale ; l'espagnol via \esp{..} / espagnol
\usepackage[a4paper,margin=2cm,top=3.4cm,bottom=2.8cm,headheight=1.4cm,footskip=1.3cm]{geometry}
\usepackage{amsmath,amssymb,amsfonts}
\usepackage{mathtools} % \xmapsto, \coloneqq... souvent utilisés par les modèles
\usepackage{cancel} % simplifications barrées (bilans de forces)
% \abs{z} (module, valeur absolue) et \norm{u} (norme), utilisés par les modèles
\providecommand{\abs}{}\let\abs\relax\DeclarePairedDelimiter{\abs}{\lvert}{\rvert}
\providecommand{\norm}{}\let\norm\relax\DeclarePairedDelimiter{\norm}{\lVert}{\rVert}
\usepackage[table]{xcolor} % [table] : fournit aussi \rowcolors (alternance de lignes)
\usepackage{pagecolor}
\usepackage{tikz}
\usetikzlibrary{arrows.meta,positioning,calc,shapes.geometric,shapes.misc,shapes.arrows,decorations.pathmorphing,decorations.pathreplacing,decorations.markings,patterns,shadows,mindmap,trees,fit,backgrounds,matrix,intersections,angles,quotes}
\usepackage{pgfplots}
\usepackage[version=4]{mhchem} % équations nucléaires \ce{^{235}_{92}U}
\usepackage[european,siunitx]{circuitikz} % schémas électriques
\pgfplotsset{compat=1.18}
\usepgfplotslibrary{fillbetween,groupplots} % aires entre courbes (calcul intégral)
\usepackage{enumitem}
\usepackage{fancyhdr}
% [most] charge toutes les bibliothèques tcolorbox (listings, minted, raster,
% documentation...) et gonfle inutilement la mémoire TeX sur un document long
% (~300 boîtes) : on ne charge que ce qui est réellement utilisé.
\usepackage[skins,breakable]{tcolorbox}
\usepackage{graphicx}
\usepackage{colortbl}
\usepackage{booktabs}
\usepackage{tabularx}
\usepackage{diagbox} % case d'en-tête diagonale des tableaux à double entrée
% Colonnes extensibles centrée / gauche / droite (C, L, R), souvent utilisées par les modèles
\newcolumntype{C}{>{\centering\arraybackslash}X}
\newcolumntype{L}{>{\raggedright\arraybackslash}X}
\newcolumntype{R}{>{\raggedleft\arraybackslash}X}
\usepackage{array}
\usepackage{multicol}
\usepackage{siunitx}
% parse-numbers=false : accepte aussi \SI{2,5 \times 10^{-3}}{...}
\sisetup{locale=FR,output-decimal-marker={,},group-minimum-digits=4,parse-numbers=false}
\DeclareSIUnit\ohm{\ensuremath{\symup{\Omega}}} % U+2126 absent de la police
\DeclareSIUnit\division{div} % réglages d'oscilloscope (ms/div, V/div)
\DeclareSIUnit\tr{tr} % tours (vitesse de rotation, ex. \SI{1500}{\tr\per\minute})
\DeclareSIUnit\molar{M} % concentration molaire (ex. \SI{3}{\molar}, \SI{1}{\micro\molar})
\DeclareSIUnit\an{an} % année (le modèle confond parfois avec \year, un compteur TeX) — ex. \SI{5}{\kilo\meter\per\an}
\DeclareSIUnit\FCFA{FCFA} % franc CFA (ex. \SI{500}{\FCFA\per\kilo\gram})
\DeclareSIUnit\degreeBrix{\textdegree Brix} % teneur en sucre d'un sirop/jus (réfractométrie)
\DeclareSIUnit\month{mois} % le modèle utilise parfois \month, un compteur TeX, comme unité de durée
\usepackage{qrcode}
% \degree : commande fréquemment utilisée par les modèles pour un angle ou une
% température en dehors de \SI{}{} (non fournie par les paquets chargés).
\providecommand{\degree}{^{\circ}}
% \qed (fin de démonstration), utilisé par les modèles sans amsthm
\providecommand{\qed}{\hfill$\square$}
% \barre : double barre des valeurs interdites dans un tableau de variations (tkz-tab)
\providecommand{\barre}{\ensuremath{\|}}
% environnement proof (amsthm non chargé) utilisé par les modèles
\ifdefined\proof\else\newenvironment{proof}[1][Démonstration]{\par\noindent\textit{#1.}\ }{\qed\par}\fi
\usepackage{lastpage}
\usepackage{hyperref}
\hypersetup{hidelinks}

% ============================================================
% Palette « Pop » OnBuch+ — mêmes hex que la classe P (plus_ui.dart)
% ============================================================
\definecolor{poporange}{HTML}{F59321}
\definecolor{popdark}{HTML}{E07A0C}
\definecolor{popcream}{HTML}{FFF6E8}
\definecolor{popink}{HTML}{1C1714}
\definecolor{poppurple}{HTML}{7A5AE0}
\definecolor{popgreen}{HTML}{2E9E6A}
\definecolor{poppink}{HTML}{E0457B}
\definecolor{popblue}{HTML}{2F7DE1}
\definecolor{popgold}{HTML}{C98A12}
\definecolor{popmuted}{HTML}{8A7F76}
\definecolor{popline}{HTML}{EADFD2}
\definecolor{popgray}{HTML}{8A7F76}
\colorlet{popgrey}{popgray}
% Alias tolérants (le modèle nomme parfois les couleurs différemment)
\colorlet{popwhite}{white}
\colorlet{popblack}{popink}
\colorlet{popred}{poppink}
\colorlet{popyellow}{popgold}
% Teintes claires (fonds de tableaux/encadrés) — le modèle nomme parfois
% les couleurs avec un suffixe L (ex. popblueL) sans les avoir définies.
\colorlet{poporangeL}{poporange!20}
\colorlet{popdarkL}{popdark!20}
\colorlet{poppurpleL}{poppurple!20}
\colorlet{popgreenL}{popgreen!20}
\colorlet{poppinkL}{poppink!20}
\colorlet{popblueL}{popblue!20}
\colorlet{popgoldL}{popgold!20}
\colorlet{popredL}{popred!20}
\colorlet{popgrayL}{popgray!20}
% Teintes claires pour fonds de tableaux / zones
\colorlet{popblueL}{popblue!10!white}
\colorlet{poporangeL}{poporange!14!white}
\colorlet{popgreenL}{popgreen!12!white}
\colorlet{poppurpleL}{poppurple!12!white}
\colorlet{poppinkL}{poppink!12!white}
\colorlet{popgoldL}{popgold!14!white}

\pagecolor{popcream}

% ============================================================
% Typographie
% ============================================================
\defaultfontfeatures{Ligatures=TeX}
\setmainfont{PlusJakartaSans-Medium.ttf}[Path=../../../fonts/,BoldFont=PlusJakartaSans-Bold.ttf]
% Maths en sans-serif (Fira Math) avec chiffres et lettres droites de Plus
% Jakarta, pour que les formules se fondent dans le texte.
\usepackage[math-style=ISO,bold-style=ISO]{unicode-math}
\setmathfont{FiraMath-Regular.otf}[Path=../../../fonts/]
\setmathfont{PlusJakartaSans-Medium.ttf}[Path=../../../fonts/,range={up/num,up/{Latin,latin},\mathup}]
\usepackage{icomma} % virgule décimale sans espace en mode math
% Caractères mathématiques Unicode absents des polices de texte (Plus Jakarta,
% Archivo Black) : rendus par leur équivalent mathématique, en texte comme en
% formule — sinon ils s'affichent en carré vide (ex. « Arithmétique dans ℤ »).
\usepackage{newunicodechar}
\newunicodechar{ℝ}{\ensuremath{\mathbb{R}}}
\newunicodechar{ℤ}{\ensuremath{\mathbb{Z}}}
\newunicodechar{ℂ}{\ensuremath{\mathbb{C}}}
\newunicodechar{ℕ}{\ensuremath{\mathbb{N}}}
\newunicodechar{ℚ}{\ensuremath{\mathbb{Q}}}
\newunicodechar{𝔻}{\ensuremath{\mathbb{D}}}
\newunicodechar{α}{\ensuremath{\alpha}}
\newunicodechar{β}{\ensuremath{\beta}}
\newunicodechar{γ}{\ensuremath{\gamma}}
\newunicodechar{δ}{\ensuremath{\delta}}
\newunicodechar{ε}{\ensuremath{\varepsilon}}
\newunicodechar{θ}{\ensuremath{\theta}}
\newunicodechar{λ}{\ensuremath{\lambda}}
\newunicodechar{μ}{\ensuremath{\mu}}
\newunicodechar{σ}{\ensuremath{\sigma}}
\newunicodechar{τ}{\ensuremath{\tau}}
\newunicodechar{φ}{\ensuremath{\varphi}}
\newunicodechar{ψ}{\ensuremath{\psi}}
\newunicodechar{ω}{\ensuremath{\omega}}
\newunicodechar{Σ}{\ensuremath{\Sigma}}
% majuscules produites par \MakeUppercase dans les titres (α-aminés -> Α)
\newunicodechar{Α}{\ensuremath{\alpha}}
\newunicodechar{Β}{\ensuremath{\beta}}
\newunicodechar{Γ}{\ensuremath{\Gamma}}
\newunicodechar{Δ}{\ensuremath{\Delta}}
\newunicodechar{Ε}{\ensuremath{\varepsilon}}
\newunicodechar{Θ}{\ensuremath{\Theta}}
\newunicodechar{Λ}{\ensuremath{\Lambda}}
\newunicodechar{Μ}{\ensuremath{\mu}}
\newunicodechar{Τ}{\ensuremath{\tau}}
\newunicodechar{Φ}{\ensuremath{\Phi}}
\newunicodechar{Ψ}{\ensuremath{\Psi}}
\newunicodechar{Ω}{\ensuremath{\Omega}}
\newunicodechar{↦}{\ensuremath{\mapsto}}
\newunicodechar{∈}{\ensuremath{\in}}
\newunicodechar{∉}{\ensuremath{\notin}}
\newunicodechar{∧}{\ensuremath{\wedge}}
\newunicodechar{⇔}{\ensuremath{\Leftrightarrow}}
\newunicodechar{⟺}{\ensuremath{\Longleftrightarrow}}
\newunicodechar{⇒}{\ensuremath{\Rightarrow}}
\newunicodechar{⟹}{\ensuremath{\Longrightarrow}}
\newunicodechar{∘}{\ensuremath{\circ}}
\newunicodechar{∪}{\ensuremath{\cup}}
\newunicodechar{∩}{\ensuremath{\cap}}
\newunicodechar{≡}{\ensuremath{\equiv}}
\newunicodechar{⊂}{\ensuremath{\subset}}
\newunicodechar{⊥}{\ensuremath{\perp}}
\newunicodechar{∃}{\ensuremath{\exists}}
\newunicodechar{∀}{\ensuremath{\forall}}
\newunicodechar{∛}{\ensuremath{\sqrt[3]{\,}}}
\newunicodechar{∜}{\ensuremath{\sqrt[4]{\,}}}
\newunicodechar{⃗}{\ensuremath{{}^{\to}}}
\newunicodechar{ⁿ}{\ifmmode^{n}\else\textsuperscript{\ensuremath{n}}\fi}
\newunicodechar{ˣ}{\ifmmode^{x}\else\textsuperscript{\ensuremath{x}}\fi}
\newunicodechar{ᵇ}{\ifmmode^{b}\else\textsuperscript{\ensuremath{b}}\fi}
\newunicodechar{ʳ}{\ifmmode^{r}\else\textsuperscript{\ensuremath{r}}\fi}
\newunicodechar{ᵘ}{\ifmmode^{u}\else\textsuperscript{\ensuremath{u}}\fi}
\newunicodechar{ʸ}{\ifmmode^{y}\else\textsuperscript{\ensuremath{y}}\fi}
\newunicodechar{ᵃ}{\ifmmode^{a}\else\textsuperscript{\ensuremath{a}}\fi}
\newunicodechar{ᵉ}{\ifmmode^{e}\else\textsuperscript{\ensuremath{e}}\fi}
\newunicodechar{ᵏ}{\ifmmode^{k}\else\textsuperscript{\ensuremath{k}}\fi}
\newunicodechar{ᵖ}{\ifmmode^{p}\else\textsuperscript{\ensuremath{p}}\fi}
\newunicodechar{ᴺ}{\ifmmode^{N}\else\textsuperscript{\ensuremath{N}}\fi}
\newunicodechar{ᶜ}{\ifmmode^{c}\else\textsuperscript{\ensuremath{c}}\fi}
\newunicodechar{ʰ}{\ifmmode^{h}\else\textsuperscript{\ensuremath{h}}\fi}
\newunicodechar{ᵠ}{\ifmmode^{q}\else\textsuperscript{\ensuremath{q}}\fi}
\newunicodechar{ₙ}{\ifmmode_{n}\else\textsubscript{\ensuremath{n}}\fi}
\newunicodechar{ₐ}{\ifmmode_{a}\else\textsubscript{\ensuremath{a}}\fi}
\newunicodechar{ᵢ}{\ifmmode_{i}\else\textsubscript{\ensuremath{i}}\fi}
\newunicodechar{ᵣ}{\ifmmode_{r}\else\textsubscript{\ensuremath{r}}\fi}
\newunicodechar{ₖ}{\ifmmode_{k}\else\textsubscript{\ensuremath{k}}\fi}
\newunicodechar{ᵦ}{\ifmmode_{\beta}\else\textsubscript{\ensuremath{\beta}}\fi}
\newunicodechar{ₓ}{\ifmmode_{x}\else\textsubscript{\ensuremath{x}}\fi}
\newfontfamily\popdisplay{ArchivoBlack-Regular.ttf}[Path=../../../fonts/]
\newfontfamily\poplogo{SpaceGrotesk-Bold.ttf}[Path=../../../fonts/]
\newfontfamily\popsemi{PlusJakartaSans-SemiBold.ttf}[Path=../../../fonts/]
\newfontfamily\popxbold{PlusJakartaSans-ExtraBold.ttf}[Path=../../../fonts/]
\newfontfamily\popxboldls{PlusJakartaSans-ExtraBold.ttf}[Path=../../../fonts/,LetterSpace=3.0]

\setlist[enumerate]{leftmargin=1.7em,itemsep=5pt,topsep=4pt}
\setlist[itemize]{leftmargin=1.7em,itemsep=4pt,topsep=4pt,label={\color{poporange}\textbullet}}
\setlength{\parindent}{0pt}
\setlength{\parskip}{5pt}
\renewcommand{\arraystretch}{1.25}

% pgfplots : style Pop par défaut
\pgfplotsset{
  every axis/.append style={axis line style={popink,line width=1pt},tick style={popink},
    grid style={popline,line width=0.5pt},label style={font=\small},tick label style={font=\footnotesize}},
  popcurve/.style={poporange,line width=1.8pt,no markers,smooth},
}
% Même style disponible dans un \draw TikZ simple (hors pgfplots)
\tikzset{popcurve/.style={poporange,line width=1.8pt}}
\tikzset{popgrid/.style={popline,very thin}}
\colorlet{popcurve}{poporange} % le nom du style est parfois utilisé comme couleur

% ============================================================
% Pill / squiggle
% ============================================================
\newcommand{\poppill}[3]{%
  \tikz[baseline=-0.55ex]\node[draw=popink,line width=1.2pt,rounded corners=2.6pt,%
    fill=#2,inner xsep=6.5pt,inner ysep=3pt]{%
    \popxboldls\fontsize{7}{7}\selectfont\color{#3}\MakeUppercase{#1}};%
}
\newcommand{\popsquiggle}[1]{%
  \tikz[baseline=-0.4ex]{
    \draw[#1,line width=2.6pt,line cap=round]
      plot [smooth] coordinates {(0,0.09) (0.34,0.01) (0.68,0.21) (1.02,0.01) (1.36,0.10)};
  }%
}
% Mot-clé surligné (marqueur orange sous le texte)
\newcommand{\cle}[1]{{\popxbold\color{popdark}#1}}

% ============================================================
% En-tête / pied
% ============================================================
\pagestyle{fancy}
\fancyhf{}
\renewcommand{\headrulewidth}{0pt}
\renewcommand{\footrulewidth}{0pt}
\lhead{\raisebox{-0.15ex}{\poplogo\fontsize{13}{13}\selectfont\color{popink}OnBuch\color{poporange}+}}
\rhead{\poppill{\DOCMATIERE\ \textbullet\ \DOCNIVEAU\ \textbullet\ Leçon \DOCLECON}{white}{popink}}
\lfoot{\popxbold\fontsize{7.6}{7.6}\selectfont\color{popmuted}\MakeUppercase{Cours signé OnBuch+}\ \textbullet\ {\popsemi\DOCTITRE}}
\rfoot{\tikz[baseline=-0.6ex]\node[rounded corners=8.5pt,fill=popink,minimum height=17pt,minimum width=17pt,inner xsep=5pt,inner sep=0]{\popxbold\fontsize{8}{8}\selectfont\color{popcream}\thepage\,/\,\pageref*{LastPage}};}

% ============================================================
% Titres
% ============================================================
\newcommand{\chaptitre}[2]{%
  \begin{tcolorbox}[enhanced,colback=poporange,colframe=popink,boxrule=2.6pt,arc=11pt,
      drop shadow=popink,left=15pt,right=15pt,top=12pt,bottom=11pt,before skip=3pt,after skip=18pt]
    \poppill{#2}{popink}{popcream}\par\vspace{9pt}
    {\raggedright\hyphenpenalty=10000\exhyphenpenalty=10000\popdisplay\fontsize{22}{24}\selectfont\color{popcream}\MakeUppercase{#1}\par}
  \end{tcolorbox}
}
% Section numérotée : pastille orange avec le numéro + titre Archivo
\newcounter{popsec}
\newcommand{\coursec}[1]{%
  \stepcounter{popsec}\setcounter{popsub}{0}%
  \par\vspace{16pt}\noindent
  \begin{minipage}{\linewidth}
  \tikz[baseline=-0.7ex]\node[draw=popink,line width=1.6pt,fill=poporange,rounded corners=4pt,minimum size=22pt,inner sep=0]{\popdisplay\fontsize{12}{12}\selectfont\color{popcream}\thepopsec};%
  \hspace{8pt}{\popdisplay\fontsize{15}{17}\selectfont\color{popink}\MakeUppercase{#1}}\par
  \vspace{4pt}\noindent{\color{poporange}\rule{36pt}{3.6pt}}
  \end{minipage}\par\nopagebreak\vspace{8pt}
}
\newcounter{popsub}
\newcommand{\courssub}[1]{%
  \stepcounter{popsub}%
  \par\vspace{9pt}\noindent{\popxbold\fontsize{11.5}{13}\selectfont\color{popink}{\color{poporange}\thepopsec.\thepopsub}\hspace{6pt}#1}\par\nopagebreak\vspace{4pt}
}

% ============================================================
% Boîtes pédagogiques — même recette : fond blanc, bordure épaisse colorée,
% ombre dure encre, badge plein. Titre optionnel [#1].
% ============================================================
\tcbset{popbase/.style={breakable,enhanced,colback=white,boxrule=2.3pt,arc=9pt,
  shadow={3pt}{-3pt}{0pt}{popink},left=14pt,right=14pt,top=11pt,bottom=11pt,before skip=11pt,after skip=15pt,
  fontupper=\popsemi\fontsize{10.6}{14.5}\selectfont\color{popink},
  fontlower=\popsemi\fontsize{10.6}{14.5}\selectfont\color{popink}}}
% Coche dessinée (le glyphe ✓ n'existe pas dans les polices du document)
\newcommand{\popcheck}[1]{\tikz[baseline=-0.5ex]\draw[#1,line width=1.6pt,line cap=round,line join=round](-0.13,0.0)--(-0.04,-0.09)--(0.13,0.1);}
\providecommand{\coche}{\popcheck{popgreen}}
\providecommand{\resultat}[1]{\par\smallskip\noindent\fcolorbox{popgreen}{popgreenL}{\parbox{\dimexpr\linewidth-2\fboxsep-2\fboxrule\relax}{\textbf{Résultat.} #1}}\par\smallskip}
\newcommand{\popboxhead}[3]{\poppill{#1}{#2}{white}\ifx\relax#3\relax\else\hspace{6pt}{\popxbold\fontsize{10.6}{12}\selectfont\color{popink}#3}\fi\par\vspace{7pt}}

\newtcolorbox{aretenir}[1][]{popbase,colframe=popgreen,before upper={\popboxhead{À retenir}{popgreen}{#1}}}
% Texte en espagnol (document de lecture, dialogue, poème, extrait) : cadre
% d'étude. Un titre optionnel (source, auteur) s'affiche dans l'en-tête.
\newtcolorbox{textebox}[1][]{popbase,colframe=popblue,before upper={\popboxhead{Texto}{popblue}{#1}\begin{otherlanguage}{spanish}},after upper={\end{otherlanguage}}}
% Marques ✓ / ✗ (forme correcte / fautive) souvent utilisées par les modèles
\providecommand{\cross}{\textcolor{poppink}{\ensuremath{\times}}}
\providecommand{\xmark}{\cross}\providecommand{\crossmark}{\cross}\providecommand{\cmark}{\ensuremath{\checkmark}}
% Ligne de réponse / justification à compléter (inventée par les modèles)
\providecommand{\justification}[1]{\par\noindent\textit{Justification~:} \dotfill\par}
% \textipa (package tipa non chargé) : la police ne couvre pas l'API -> simple passage
\providecommand{\textipa}[1]{#1}
% Mot ou phrase en espagnol à l'intérieur d'un paragraphe français
\newcommand{\esp}[1]{\begingroup\selectlanguage{spanish}\itshape #1\endgroup}
\newtcolorbox{exemplebox}[1][]{popbase,colframe=poporange,before upper={\popboxhead{Exemple}{poporange}{#1}}}
\newtcolorbox{definition}[1][]{popbase,colframe=popblue,before upper={\popboxhead{Définition}{popblue}{#1}}}
\newtcolorbox{propriete}[1][]{popbase,colframe=poppurple,before upper={\popboxhead{Propriété}{poppurple}{#1}}}
\newtcolorbox{methode}[1][]{popbase,colframe=popgold,before upper={\popboxhead{Méthode}{popgold}{#1}}}
\newtcolorbox{attention}[1][]{popbase,colframe=poppink,before upper={\popboxhead{Attention}{poppink}{#1}}}
\newtcolorbox{experience}[1][]{popbase,colframe=popblue,colback=popblueL,before upper={\popboxhead{Expérience}{popblue}{#1}}}
% « Le savais-tu ? » : carte violette pleine (contexte, culture, Cameroun)
\newtcolorbox{savaistu}[1][]{popbase,colback=poppurple,colframe=popink,
  fontupper=\popsemi\fontsize{10.4}{14.2}\selectfont\color{white},
  before upper={\poppill{Le savais-tu\,?}{white}{poppurple}\ifx\relax#1\relax\else\hspace{6pt}{\popxbold\color{white}#1}\fi\par\vspace{7pt}}}
% Exercice résolu : énoncé en haut, \tcblower puis la solution sur fond crème
\newcounter{popexo}\newcounter{popexr}
\newtcolorbox{exoresolu}[1][]{popbase,colframe=poporange,colbacklower=popcream,
  segmentation style={popink,line width=1.2pt,dashed},
  before upper={\stepcounter{popexr}\popboxhead{Exercice résolu \thepopexr}{poporange}{#1}},
  before lower={\poppill{Solution}{popink}{popcream}\par\vspace{6pt}}}
% Exercice (non résolu en place) — la correction est dans la partie Corrigés
\newtcolorbox{exercice}[1][]{popbase,colframe=popink,
  before upper={\stepcounter{popexo}\popboxhead{Exercice \thepopexo}{popink}{#1}}}
% Corrigé
\newtcolorbox{corrige}[1][]{popbase,colframe=popgreen,colback=popgreenL,
  before upper={\popboxhead{Corrigé}{popgreen}{#1}}}
% Objectifs / prérequis / bilan
\newtcolorbox{objectifs}{popbase,colframe=popink,colback=poporangeL,
  before upper={\popboxhead{Objectifs de la leçon}{popink}{}}}
\newtcolorbox{prerequis}{popbase,colframe=popink,colback=popblueL,
  before upper={\popboxhead{Prérequis}{popblue}{}}}
\newtcolorbox{bilan}{popbase,colframe=popink,colback=white,boxrule=2.8pt,
  before upper={\popboxhead{Fiche bilan}{poporange}{L'essentiel à connaître pour l'examen}}}
% Figure encadrée avec légende
\newtcolorbox{popfigure}[1][]{enhanced,breakable=false,colback=white,colframe=popline,boxrule=1.4pt,arc=8pt,
  left=10pt,right=10pt,top=10pt,bottom=8pt,before skip=10pt,after skip=12pt,halign=center,
  lower separated=false,halign lower=center,
  fontlower=\popsemi\fontsize{9}{11.5}\selectfont\color{popmuted},
  #1}
\newcommand{\legende}[1]{\par\vspace{6pt}{\popsemi\fontsize{9}{11.5}\selectfont\color{popmuted}#1}}

% ============================================================
% Signature OnBuch+
% ============================================================
\newcommand{\poplogoinline}[1]{{\poplogo\fontsize{#1}{#1}\selectfont\color{popink}OnBuch\color{poporange}+}}
\newcommand{\signatureonbuch}{%
  \begin{tcolorbox}[enhanced,colback=popink,colframe=popink,boxrule=2.2pt,arc=10pt,
      drop shadow=poporange,left=14pt,right=14pt,top=10pt,bottom=10pt,before skip=0pt,after skip=0pt]
    \tikz[baseline=-0.7ex]\node[circle,fill=popgreen,draw=popcream,line width=1.3pt,minimum size=22pt,inner sep=0]{\popcheck{white}};%
    \hspace{9pt}\begin{minipage}[c]{0.62\linewidth}
      {\popxbold\fontsize{12}{13}\selectfont\color{popcream}Signé\ \textbullet\ {\poplogo OnBuch\color{poporange}+}}\par\vspace{2pt}
      {\raggedright\popsemi\fontsize{8}{10.5}\selectfont\color{popline}\MakeUppercase{Équipe pédagogique OnBuch+}\\\MakeUppercase{Conforme au programme officiel MINESEC}\par}
    \end{minipage}\hfill
    {\poplogo\fontsize{22}{22}\selectfont\color{popcream}OnBuch\color{poporange}+}
  \end{tcolorbox}%
}

% ============================================================
% Page de garde
% #1 titre · #2 matière · #3 niveau · #4 module · #5 n° leçon · #6 durée ·
% #7 description / objectifs (texte ou itemize)
% ============================================================
\newcommand{\pagegarde}[7]{%
  \thispagestyle{empty}
  \newgeometry{a4paper,margin=1.7cm,top=1.6cm,bottom=1.4cm}%
  \begin{tikzpicture}[remember picture,overlay]
    \fill[poporange!18] ([xshift=-3.2cm,yshift=-3.6cm]current page.north east) circle (4.6cm);
    \fill[poppurple!14] ([xshift=2.4cm,yshift=7.6cm]current page.south west) circle (3.8cm);
  \end{tikzpicture}%
  {\poplogo\fontsize{30}{30}\selectfont\color{popink}OnBuch\color{poporange}+}\hfill
  \raisebox{0.6ex}{\poppill{Cours signé \textbullet\ Premium}{popink}{popcream}}\par
  \vspace{1pt}\popsquiggle{poppurple}\par
  \vspace{8pt}
  \begin{tcolorbox}[enhanced,colback=poporange,colframe=popink,boxrule=2.8pt,arc=13pt,
      drop shadow=popink,left=18pt,right=18pt,top=12pt,bottom=13pt,after skip=10pt]
    \poppill{#2 \textbullet\ #3}{popink}{popcream}\hspace{5pt}\poppill{#4}{white}{popink}\par\vspace{8pt}
    {\popdisplay\fontsize{14}{14}\selectfont\color{popink}LEÇON #5}\par\vspace{4pt}
    {\raggedright\hyphenpenalty=10000\exhyphenpenalty=10000\popdisplay\fontsize{25}{27}\selectfont\color{popcream}\MakeUppercase{#1}\par}
    \vspace{8pt}
    {\popxbold\fontsize{9.5}{9.5}\selectfont\color{popink}\MakeUppercase{Durée conseillée : #6}}
  \end{tcolorbox}
  \begin{tcolorbox}[enhanced,colback=white,colframe=popink,boxrule=2.2pt,arc=10pt,
      drop shadow=popink,left=16pt,right=16pt,top=10pt,bottom=10pt,after skip=8pt]
    \poppill{Dans cette leçon}{poppurple}{white}\par\vspace{5pt}
    \popsemi\fontsize{9.3}{12.6}\selectfont\color{popink}#7
  \end{tcolorbox}
  \noindent\begin{minipage}[t]{0.487\linewidth}
    \begin{tcolorbox}[enhanced,colback=popgreen,colframe=popink,boxrule=2.2pt,arc=10pt,
        drop shadow=popink,left=11pt,right=11pt,top=10pt,bottom=10pt,height=3.1cm,valign=center]
      \tcbox[colback=white,colframe=popink,boxrule=1.3pt,arc=5pt,boxsep=0pt,nobeforeafter,
        left=3pt,right=3pt,top=3pt,bottom=3pt,valign=center]{\qrcode[height=1.9cm]{https://play.google.com/store/apps/details?id=com.luvvix.onbuch&hl=fr}}%
      \hspace{8pt}\begin{minipage}[c]{0.5\linewidth}
        {\popdisplay\fontsize{11}{12.5}\selectfont\color{white}\MakeUppercase{Télécharge OnBuch}}\par\vspace{4pt}
        {\raggedright\popsemi\fontsize{8.4}{11}\selectfont\color{white}Gratuit \textbullet\ Google Play\\Tuteur IA, annales, concours, résultats\par}
      \end{minipage}
    \end{tcolorbox}
  \end{minipage}\hfill
  \begin{minipage}[t]{0.487\linewidth}
    \begin{tcolorbox}[enhanced,colback=poppurple,colframe=popink,boxrule=2.2pt,arc=10pt,
        drop shadow=popink,left=11pt,right=11pt,top=10pt,bottom=10pt,height=3.1cm,valign=center]
      \tikz[baseline=-0.7ex]\node[circle,fill=white,draw=popink,line width=1.3pt,minimum size=1.6cm,inner sep=0]{\popdisplay\fontsize{24}{24}\selectfont\color{poppurple}+};%
      \hspace{8pt}\begin{minipage}[c]{0.6\linewidth}
        {\popdisplay\fontsize{11}{12.5}\selectfont\color{white}\MakeUppercase{Passe à OnBuch+}}\par\vspace{4pt}
        {\raggedright\popsemi\fontsize{8.4}{11}\selectfont\color{white}Tous les cours signés, Tuteur IA illimité, fiches PDF — onglet OnBuch+ de l'app\par}
      \end{minipage}
    \end{tcolorbox}
  \end{minipage}
  \vspace{8pt}
  \popcontenu
  \vfill
  \signatureonbuch
  \restoregeometry
  \newpage
}

% Rangée « Ce cours contient » (page de garde)
\newcommand{\poptile}[3]{% #1 couleur #2 gros texte #3 légende
  \begin{tcolorbox}[enhanced,colback=#1,colframe=popink,boxrule=2pt,arc=9pt,shadow={2.5pt}{-2.5pt}{0pt}{popink},
      left=6pt,right=6pt,top=9pt,bottom=9pt,halign=center,valign=center,height=1.75cm,width=\linewidth,nobeforeafter]
    {\popdisplay\fontsize{12.5}{13}\selectfont\color{white}\MakeUppercase{#2}}\par\vspace{3pt}
    {\centering\popsemi\fontsize{7.8}{9.5}\selectfont\color{white}#3\par}
  \end{tcolorbox}}
\newcommand{\popcontenu}{%
  \par\noindent{\popxboldls\fontsize{7.5}{7.5}\selectfont\color{popmuted}CE COURS SIGNÉ CONTIENT}\par\vspace{7pt}
  \noindent\begin{minipage}[t]{0.235\linewidth}\poptile{popblue}{Cours}{complet, illustré, pas à pas}\end{minipage}\hfill
  \begin{minipage}[t]{0.235\linewidth}\poptile{poporange}{Méthodes}{et exercices résolus}\end{minipage}\hfill
  \begin{minipage}[t]{0.235\linewidth}\poptile{poppink}{Exercices}{type examen + corrigés}\end{minipage}\hfill
  \begin{minipage}[t]{0.235\linewidth}\poptile{popgreen}{Fiche bilan}{l'essentiel à retenir}\end{minipage}\par}

% Sommaire
\newtcolorbox{sommaire}{enhanced,colback=white,colframe=popink,boxrule=2.2pt,arc=10pt,
  drop shadow=popink,left=18pt,right=18pt,top=13pt,bottom=13pt,before skip=6pt,after skip=20pt,
  before upper={\poppill{Sommaire}{poppurple}{white}\par\vspace{9pt}\popsemi\fontsize{10.6}{14.5}\selectfont\color{popink}}}

% Signature de fin de cours — poussée tout en bas de la dernière page
\newcommand{\findecours}{%
  \par\vfill\nopagebreak
  \begin{center}\popsquiggle{poporange}\end{center}\vspace{4pt}
  \signatureonbuch
}
\AtBeginDocument{\def\llbracket{\mathopen{[\mkern-3.5mu[}}\def\rrbracket{\mathclose{]\mkern-3.5mu]}}} % intervalles d'entiers (glyphe ⟦ absent de la police)
% Glyphes absents de Fira Math : redéfinis à partir de glyphes disponibles
\AtBeginDocument{%
  \def\setminus{\mathbin{\backslash}}\let\smallsetminus\setminus
  \def\vdots{\mathord{\vbox{\baselineskip3.2pt\lineskiplimit0pt\kern4pt\hbox{.}\hbox{.}\hbox{.}}}}%
  \def\ddots{\mathinner{\mkern1mu\raise6.5pt\hbox{.}\mkern2mu\raise3.6pt\hbox{.}\mkern2mu\raise0.7pt\hbox{.}\mkern1mu}}%
  \def\triangle{\mathord{\scalebox{0.9}{$\symup{\Delta}$}}}%
  \def\nabla{\mathord{\raisebox{\depth}{\scalebox{1}[-1]{$\symup{\Delta}$}}}}%
  \def\checkmark{\popcheck{popdark}}%
  \def\star{\mathbin{\ast}}\def\bigstar{\mathord{\ast}}%
  \def\diamond{\mathbin{\scalebox{0.6}{$\Diamond$}}}%
  \def\bigcup{\mathop{\mathchoice{\raisebox{-0.3em}{\scalebox{1.7}{$\cup$}}}{\scalebox{1.25}{$\cup$}}{\cup}{\cup}}\displaylimits}%
  \def\bigcap{\mathop{\mathchoice{\raisebox{-0.3em}{\scalebox{1.7}{$\cap$}}}{\scalebox{1.25}{$\cap$}}{\cap}{\cap}}\displaylimits}%
  \def\lesseqqgtr{\mathrel{\lessgtr}}\def\bigoplus{\mathop{\oplus}\displaylimits}%
}
\providecommand{\ssi}{si et seulement si} % abréviation fréquente
\AtBeginDocument{\providecommand{\wideparen}{\overparen}} % arc de cercle

\begin{document}

\pagegarde{Lectura y comentario : relaciones internacionales, relaciones Norte-Sur, sistemas políticos, derechos y deberes}{Espagnol}{Tle A}{Módulo 3 : Vie citoyenne et ouverture au monde}{E13}{2 h}{Cette leçon propose un texte authentique sur les relations internationales et les systèmes politiques, suivi d'activités de compréhension, de lexique thématique, de méthode de commentaire et de production écrite guidée, conformes au programme de Terminale A (2 h).
\vspace{4pt}
\begin{multicols}{2}\small
\begin{itemize}
\item Lire et comprendre un texte d'actualité politique en espagnol.
\item Identifier et expliquer le vocabulaire des relations internationales (ONU, UA, globalisation, dette, développement).
\item Différencier les principaux systèmes politiques : démocratie, dictature, monarchie, république.
\item Citer les droits, devoirs et libertés fondamentaux du citoyen dans un régime démocratique.
\item Appliquer la méthode du commentaire structuré (introduction, développement, conclusion).
\item Rédiger un commentaire argumenté de 150‑200 mots en espagnol.
\end{itemize}\end{multicols}}

\begin{sommaire}
\begin{multicols}{2}
\begin{enumerate}
\item 1. Texte support \& glossaire
\item 2. Compréhension détaillée
\item 3. Lexique thématique par champs
\item 4. Méthode du commentaire \& commentaire modèle
\item 5. Production guidée : prise de position citoyenne
\item Activité d'intégration
\item Méthodes et exercices résolus
\item Exercices
\item Corrigés des exercices
\item Fiche bilan
\end{enumerate}
\end{multicols}
\end{sommaire}

\begin{objectifs}
\begin{itemize}
\item Lire et comprendre un texte d'actualité politique en espagnol.
\item Identifier et expliquer le vocabulaire des relations internationales (ONU, UA, globalisation, dette, développement).
\item Différencier les principaux systèmes politiques : démocratie, dictature, monarchie, république.
\item Citer les droits, devoirs et libertés fondamentaux du citoyen dans un régime démocratique.
\item Appliquer la méthode du commentaire structuré (introduction, développement, conclusion).
\item Rédiger un commentaire argumenté de 150‑200 mots en espagnol.
\item Produire une courte prise de position écrite (article, lettre) sur un thème Nord‑Sud.
\end{itemize}
\end{objectifs}

\begin{prerequis}
\begin{itemize}
\item Connaître les temps du passé (pretérito indefinido, imperfecto) pour situer les événements historiques.
\item Maîtriser les connecteurs logiques (por tanto, sin embargo, además) pour structurer un commentaire.
\item Avoir étudié en Première les notions de citoyenneté et de droits de l'homme.
\item Savoir utiliser un dictionnaire bilingue pour le lexique spécialisé.
\item Disposer de bases sur la géopolitique mondiale (blocs Nord/Sud, organisations internationales).
\end{itemize}
\end{prerequis}

\newpage

\coursec{Texte support \& glossaire}

\courssub{Situation d'accroche et problématique}

Ce soir, à Yaoundé, la famille Ngo Bassa regarde le journal télévisé : on y parle du sommet de l'\esp{Unión Africana} et d'une réunion de l'\esp{ONU} sur la dette des pays africains. Le lendemain, au lycée, le professeur d'espagnol distribue un article de presse hispanophone qui traite exactement des mêmes sujets : la \esp{cooperación Norte-Sur}, la \esp{deuda externa}, la \esp{democracia}. Car ces débats ne sont pas seulement camerounais : en Espagne et en Amérique latine, les mêmes enjeux — coopération Nord-Sud, dette extérieure, démocratie — animent les journaux et les réseaux sociaux. Comprendre ces textes permet de participer aux conversations citoyennes, du lycée de Yaoundé aux universités de Madrid ou de Mexico.

\textbf{Problématique :} comment lire et comprendre un texte d'opinion en espagnol sur les relations internationales, et comment s'approprier le vocabulaire qui permet d'en parler avec précision ?

\begin{methode}[Lire efficacement un texte politique]
Avant même de lire le texte en détail, adopte la démarche suivante :
\begin{enumerate}
\item \textbf{Titre et source} : identifier le titre, le type de document (article d'opinion, éditorial, reportage), la source et la date.
\item \textbf{Lecture globale} : lire une première fois en silence, sans dictionnaire, pour saisir l'idée générale.
\item \textbf{Repérage des idées-forces} : souligner les mots-clés et repérer la structure (introduction, développement, conclusion).
\item \textbf{Annotation du vocabulaire} : entourer les termes inconnus, deviner leur sens grâce au contexte, puis vérifier dans le glossaire.
\end{enumerate}
\end{methode}

\courssub{Présentation du texte support}

Le texte que nous allons étudier est un \cle{article d'opinion} (\esp{artículo de opinión}), c'est-à-dire un texte dans lequel un journaliste ou un expert prend position et défend une thèse. Il est publié dans un journal hispanophone, dans la rubrique « Opinión ». Son registre est \cle{formel} et \cle{argumentatif} : l'auteur utilise des connecteurs logiques (\esp{sin embargo}, \esp{por lo tanto}, \esp{además}), un vocabulaire soutenu et des exemples concrets pour convaincre le lecteur.

\textbf{Première lecture (silencieuse)} : lis le texte une fois sans t'arrêter sur les mots inconnus. \textbf{Deuxième lecture (à haute voix)} : lis à voix haute en soignant la prononciation — attention notamment à \esp{globalización} (le \emph{z} se prononce comme le \emph{th} anglais en Espagne, comme un \emph{s} en Amérique latine), à \esp{cooperación} (accent tonique sur la dernière syllabe, d'où l'accent écrit) et à \esp{desarrollo} (le \emph{rr} roulé).

\begin{textebox}[La cooperación Norte-Sur ante los retos de la globalización — artículo de opinión]
En las últimas décadas, la globalización ha transformado profundamente las relaciones internacionales. Los intercambios comerciales, culturales y tecnológicos se han multiplicado, pero los beneficios de este proceso no se reparten de manera justa: mientras los países del Norte acumulan riqueza y poder, muchos países del Sur siguen enfrentándose a la pobreza, a la deuda externa y a la falta de infraestructuras.

La cooperación Norte-Sur aparece así como una necesidad urgente. Organismos internacionales como la ONU promueven programas de desarrollo y de ayuda humanitaria. Sin embargo, numerosos analistas critican una ayuda que a veces llega con condiciones políticas o económicas que limitan la soberanía de los países beneficiarios. En África, la Unión Africana defiende una cooperación basada en el respeto mutuo y en la solidaridad, y no en la dependencia.

Además, la globalización plantea nuevos retos: el cambio climático, las migraciones y las crisis sanitarias no conocen fronteras. Ningún país, por poderoso que sea, puede resolver estos problemas solo. Por lo tanto, es indispensable construir un sistema internacional más justo, donde los derechos y los deberes sean compartidos por todos.

En conclusión, la cooperación entre el Norte y el Sur no debe ser una simple caridad, sino una verdadera alianza entre socios iguales. Los ciudadanos, especialmente los jóvenes, tienen un papel fundamental: informarse, debatir y exigir a sus gobiernos políticas solidarias y democráticas. El futuro de las relaciones internacionales depende también de nosotros.
\end{textebox}

\textbf{Traduction du texte (pour vérification) :}

« Au cours des dernières décennies, la mondialisation a profondément transformé les relations internationales. Les échanges commerciaux, culturels et technologiques se sont multipliés, mais les bénéfices de ce processus ne sont pas répartis de manière juste : tandis que les pays du Nord accumulent richesse et pouvoir, beaucoup de pays du Sud continuent d'affronter la pauvreté, la dette extérieure et le manque d'infrastructures.

La coopération Nord-Sud apparaît ainsi comme une nécessité urgente. Des organismes internationaux comme l'ONU promeuvent des programmes de développement et d'aide humanitaire. Cependant, de nombreux analystes critiquent une aide qui arrive parfois avec des conditions politiques ou économiques qui limitent la souveraineté des pays bénéficiaires. En Afrique, l'Union africaine défend une coopération fondée sur le respect mutuel et la solidarité, et non sur la dépendance.

De plus, la mondialisation pose de nouveaux défis : le changement climatique, les migrations et les crises sanitaires ne connaissent pas de frontières. Aucun pays, si puissant soit-il, ne peut résoudre ces problèmes seul. Par conséquent, il est indispensable de construire un système international plus juste, où les droits et les devoirs soient partagés par tous.

En conclusion, la coopération entre le Nord et le Sud ne doit pas être une simple charité, mais une véritable alliance entre des partenaires égaux. Les citoyens, en particulier les jeunes, ont un rôle fondamental : s'informer, débattre et exiger de leurs gouvernements des politiques solidaires et démocratiques. L'avenir des relations internationales dépend aussi de nous. »

\begin{attention}[Deux tournures remarquables du texte]
\begin{itemize}
\item \esp{por poderoso que sea} : tournure concessive avec \esp{por} + adjectif + \esp{que} + \textbf{subjonctif} = « si puissant soit-il », « quelque puissant qu'il soit ». Autre exemple : \esp{por difícil que sea el problema, hay que actuar.} « Si difficile que soit le problème, il faut agir. »
\item \esp{donde los derechos y los deberes sean compartidos} : subjonctif après \esp{donde} car il s'agit d'un système \emph{irréel, à construire} (il n'existe pas encore). Compare : \esp{el sistema donde los derechos \textbf{son} compartidos} (indicatif : le système existe).
\end{itemize}
\end{attention}

\courssub{Glossaire intégré : 15 termes clés}

Chaque terme est défini en espagnol (définition simple, niveau lycée) puis traduit en français. Apprends la définition espagnole : c'est elle qui te servira au commentaire de texte.

\begin{center}
\begin{tabularx}{\linewidth}{|l|X|X|}
\hline
\rowcolor{popblueL} \textbf{Término} & \textbf{Definición en español} & \textbf{Traduction française} \\
\hline
\esp{relaciones internacionales} & vínculos políticos, económicos y culturales entre los países & relations internationales \\
\hline
\esp{cooperación} & acción de trabajar juntos para alcanzar un objetivo común & coopération \\
\hline
\esp{globalización} & proceso de integración económica y cultural del mundo & mondialisation \\
\hline
\esp{desarrollo} & crecimiento económico y mejora de las condiciones de vida & développement \\
\hline
\esp{deuda externa} & dinero que un país debe a otros países u organismos & dette extérieure \\
\hline
\esp{ayuda humanitaria} & asistencia dada a poblaciones en situación de crisis & aide humanitaire \\
\hline
\esp{soberanía} & poder de un Estado de decidir por sí mismo & souveraineté \\
\hline
\esp{intercambios} & acción de dar y recibir bienes, ideas o servicios & échanges \\
\hline
\esp{pobreza} & situación de falta de recursos económicos & pauvreté \\
\hline
\esp{reto} & desafío, objetivo difícil de alcanzar & défi \\
\hline
\esp{democracia} & sistema político donde el pueblo elige a sus gobernantes & démocratie \\
\hline
\esp{derechos} & libertades y garantías reconocidas a los ciudadanos & droits \\
\hline
\esp{deberes} & obligaciones que tienen los ciudadanos & devoirs \\
\hline
\esp{solidaridad} & apoyo entre personas o países & solidarité \\
\hline
\esp{ciudadano} & persona que pertenece a un Estado y tiene derechos y deberes & citoyen \\
\hline
\end{tabularx}
\end{center}

\begin{definition}[ONU]
\esp{ONU} = \esp{Organización de las Naciones Unidas}, organisme international créé en 1945 pour maintenir la paix et la sécurité dans le monde. En français : Organisation des Nations unies. Attention au genre : on dit \esp{la ONU} (féminin, car \esp{organización} est féminin). Exemple : \esp{La ONU promueve programas de desarrollo.} « L'ONU promeut des programmes de développement. »
\end{definition}

\begin{savaistu}[L'Union africaine]
L'\esp{Unión Africana} (UA) a été créée en 2002 à Durban, en Afrique du Sud, succédant à l'OUA (Organisation de l'unité africaine, fondée en 1963 à Addis-Abeba). Elle regroupe les États du continent africain et œuvre pour la paix, l'intégration économique et la démocratie. Le Cameroun en est membre, tout comme la Guinée équatoriale, seul pays hispanophone d'Afrique subsaharienne — un pont naturel entre le monde hispanique et l'Afrique.
\end{savaistu}

\courssub{Repérage des marqueurs de structure}

Un article d'opinion bien construit suit un plan en trois parties. Repérer les \cle{marqueurs de structure} (\esp{marcadores}) permet de comprendre l'organisation du texte et de le résumer facilement.

\begin{center}
\begin{tabularx}{\linewidth}{|l|X|X|}
\hline
\rowcolor{popblueL} \textbf{Partie} & \textbf{Marqueurs dans le texte} & \textbf{Fonction} \\
\hline
Introduction & \esp{En las últimas décadas...} & présenter le sujet et le contexte \\
\hline
Développement & \esp{Sin embargo}, \esp{Además}, \esp{Por lo tanto} & arguments, nuances, exemples \\
\hline
Conclusion & \esp{En conclusión} & synthèse et ouverture \\
\hline
\end{tabularx}
\end{center}

\begin{exemplebox}[Connecteurs argumentatifs à retenir]
\begin{itemize}
\item \esp{Sin embargo, la ayuda llega con condiciones.} « Cependant, l'aide arrive avec des conditions. » (opposition)
\item \esp{Además, la globalización plantea nuevos retos.} « De plus, la mondialisation pose de nouveaux défis. » (addition)
\item \esp{Por lo tanto, es indispensable cooperar.} « Par conséquent, il est indispensable de coopérer. » (conséquence)
\item \esp{En conclusión, la cooperación debe ser una alianza.} « En conclusion, la coopération doit être une alliance. » (synthèse)
\end{itemize}
\end{exemplebox}

\begin{exemplebox}[Exemples traduits : le vocabulaire en contexte]
\begin{itemize}
\item \esp{La globalización \textbf{acelera} los intercambios.} → La mondialisation \emph{accélère} les échanges.
\item \esp{Los países del Sur \textbf{se enfrentan a} la pobreza.} → Les pays du Sud \emph{affrontent} la pauvreté.
\item \esp{La ONU \textbf{promueve} programas de desarrollo.} → L'ONU \emph{promeut} des programmes de développement.
\item \esp{Los ciudadanos \textbf{exigen} políticas solidarias.} → Les citoyens \emph{exigent} des politiques solidaires.
\item \esp{Ningún país puede resolver estos problemas \textbf{solo}.} → Aucun pays ne peut résoudre ces problèmes \emph{seul}.
\end{itemize}
\end{exemplebox}

\begin{propriete}[Le verbe \esp{enfrentarse a}]
\esp{Enfrentarse a} + nom = « affronter, faire face à ». C'est un verbe \textbf{pronominal} : le pronom réfléchi change avec le sujet. Présent : \esp{me enfrento}, \esp{te enfrentas}, \esp{se enfrenta}, \esp{nos enfrentamos}, \esp{os enfrentáis}, \esp{se enfrentan}. Exemple : \esp{Nos enfrentamos a muchos retos.} « Nous affrontons de nombreux défis. » Ne pas oublier la préposition \esp{a} : on dit \esp{se enfrentan \textbf{a} la pobreza}.
\end{propriete}

\courssub{Analyse rapide du registre}

Le registre du texte est \cle{formel} et \cle{argumentatif}. Voici les indices qui le prouvent :

\begin{itemize}
\item \textbf{Vocabulaire soutenu} : \esp{organismos internacionales}, \esp{soberanía}, \esp{beneficiarios}, \esp{indispensable}.
\item \textbf{Connecteurs logiques variés} : \esp{sin embargo}, \esp{además}, \esp{por lo tanto}, \esp{en conclusión}.
\item \textbf{Structures impersonnelles} : \esp{es indispensable construir...}, \esp{los beneficios no se reparten...} — l'auteur généralise au lieu de dire « je ».
\item \textbf{Prise de position explicite} : \esp{no debe ser una simple caridad, sino una verdadera alianza} — l'auteur défend une thèse.
\item \textbf{Interpellation finale du lecteur} : \esp{depende también de nosotros} — procédé typique de l'article d'opinion pour impliquer le public.
\end{itemize}

\begin{attention}[Faux amis et pièges de traduction]
\begin{itemize}
\item Ne pas confondre \esp{deuda externa} (\textbf{dette extérieure}, argent dû à l'étranger) et \esp{déficit público} (\textbf{déficit public}, quand l'État dépense plus qu'il ne gagne). Les deux notions sont liées mais différentes.
\item \esp{actual} signifie « actuel, présent » et \esp{actualmente} « actuellement, de nos jours » ; en revanche, « actuel » au sens de « réel, véritable » se dit \esp{real} ou \esp{verdadero}. Exemple : \esp{la situación actual} « la situation actuelle », mais \esp{el peligro real} « le danger réel ».
\item \esp{asistencia} peut signifier « aide » (\esp{asistencia humanitaria}) mais aussi « présence, assistance » à une réunion : le contexte décide.
\item \esp{demandar} = « exiger, réclamer » (ou « poursuivre en justice »), pas « demander » poliment : pour « demander », on dit \esp{pedir}. Compare : \esp{Los ciudadanos demandan justicia.} « Les citoyens réclament justice. » / \esp{Pido información.} « Je demande des informations. »
\item \esp{exigir} (employé dans le texte) est le vrai équivalent de « exiger » : \esp{exigir a sus gobiernos políticas solidarias} « exiger de leurs gouvernements des politiques solidaires ».
\end{itemize}
\end{attention}

\begin{popfigure}
\begin{tikzpicture}[node distance=1cm]
  \node[draw=popblue,fill=popblueL,rounded corners,align=center,text width=6cm] (txt) {\textbf{Texte support}\\« La cooperación Norte-Sur ante los retos de la globalización »};
  \node[draw=poporange,fill=poporangeL,rounded corners,below=1cm of txt,align=center,text width=6cm] (glos) {\textbf{Glossaire}\\15 términos clave};
  \node[draw=popgreen,fill=popgreenL,rounded corners,below=1cm of glos,align=center,text width=6cm] (est) {\textbf{Estructura}\\introducción — desarrollo — conclusión};
  \draw[-{Stealth},thick,popblue] (txt) -- (glos);
  \draw[-{Stealth},thick,poporange] (glos) -- (est);
\end{tikzpicture}
\legende{Du texte au glossaire puis à la structure : la démarche de lecture.}
\end{popfigure}

\begin{exoresolu}[Identifier le type de texte et sa thèse]
\textbf{Énoncé.} Réponds en français, en t'appuyant sur le texte : (1) Quel est le type de ce texte et où paraît-il ? (2) Quelle thèse l'auteur défend-il ? (3) Cite en espagnol deux connecteurs qui structurent l'argumentation.
\tcblower
\textbf{Solution détaillée.}
\begin{enumerate}
\item Il s'agit d'un \textbf{article d'opinion} (\esp{artículo de opinión}), publié dans la rubrique « Opinión » d'un journal hispanophone. On le reconnaît à la prise de position de l'auteur et au registre formel et argumentatif.
\item L'auteur défend la thèse suivante : la coopération Nord-Sud ne doit pas être une simple charité (\esp{una simple caridad}) mais une véritable alliance entre partenaires égaux (\esp{una verdadera alianza entre socios iguales}), car les défis de la mondialisation (climat, migrations, crises sanitaires) exigent une solidarité partagée.
\item Deux connecteurs : \esp{sin embargo} (opposition, deuxième paragraphe) et \esp{por lo tanto} (conséquence, troisième paragraphe). On pouvait aussi citer \esp{además} et \esp{en conclusión}.
\end{enumerate}
\end{exoresolu}

\begin{exercice}[À toi de jouer]
\begin{enumerate}
\item Retrouve dans le texte un synonyme de \esp{desafío} et une expression qui désigne les « pays riches ».
\item Traduis en français : \esp{La Unión Africana defiende una cooperación basada en el respeto mutuo.}
\item Classe les mots suivants dans deux colonnes « Norte » et « Sur » selon le texte : \esp{riqueza}, \esp{pobreza}, \esp{deuda externa}, \esp{poder}, \esp{falta de infraestructuras}.
\item Relève dans le texte un nom de la famille de \esp{cooperar} et un nom de la famille de \esp{desarrollar}.
\item Traduis en espagnol : « Les jeunes citoyens doivent s'informer et débattre. »
\end{enumerate}
\end{exercice}

\begin{corrige}[Exercice « À toi de jouer »]
\begin{enumerate}
\item Synonyme de \esp{desafío} : \esp{reto}. Expression pour « pays riches » : \esp{los países del Norte} (qui \esp{acumulan riqueza y poder}).
\item « L'Union africaine défend une coopération fondée sur le respect mutuel. »
\item Norte : \esp{riqueza}, \esp{poder}. Sur : \esp{pobreza}, \esp{deuda externa}, \esp{falta de infraestructuras}.
\item \esp{cooperar} → \esp{cooperación} ; \esp{desarrollar} → \esp{desarrollo} (\esp{programas de desarrollo}).
\item \esp{Los jóvenes ciudadanos deben informarse y debatir.} (On accepte aussi \esp{Los ciudadanos jóvenes tienen que informarse y debatir.})
\end{enumerate}
\end{corrige}

\begin{aretenir}[L'essentiel de cette première étape]
\begin{itemize}
\item Un \esp{artículo de opinión} est un texte formel et argumentatif où l'auteur défend une thèse ; il se lit en quatre temps : titre et source, lecture globale, repérage des idées-forces, annotation du vocabulaire.
\item La structure type est : \esp{introducción} (\esp{en las últimas décadas...}) → \esp{desarrollo} (\esp{sin embargo}, \esp{además}, \esp{por lo tanto}) → \esp{conclusión} (\esp{en conclusión}).
\item Le lexique des relations internationales (\esp{cooperación}, \esp{desarrollo}, \esp{deuda externa}, \esp{soberanía}, \esp{globalización}) est indispensable pour comprendre et commenter l'actualité hispanophone.
\item Vigilance sur les faux amis : \esp{deuda externa} n'est pas \esp{déficit público}, \esp{demandar} ne veut pas dire « demander », et \esp{actual} ne veut pas dire « réel ».
\item Deux tournures de niveau Terminale à réutiliser au commentaire : \esp{por} + adjectif + \esp{que} + subjonctif (\esp{por poderoso que sea}) et le subjonctif après \esp{donde} pour un monde à construire (\esp{donde los derechos sean compartidos}).
\end{itemize}
\end{aretenir}

\coursec{Compréhension détaillée}

Dans la section précédente, nous avons lu le texte support \esp{« Las relaciones Norte-Sur : entre cooperación y desigualdad »}, dégagé son thème général et construit son glossaire (\esp{cooperación}, \esp{deuda}, \esp{desarrollo}, \esp{globalización}...). Nous passons maintenant à l'étape décisive de toute épreuve de commentaire au baccalauréat : la \cle{compréhension détaillée}. Il ne s'agit plus de survoler le texte, mais de le questionner ligne par ligne, de vérifier que chaque information, chaque chiffre, chaque connecteur a été compris. C'est cette lecture fine qui nourrira ensuite le commentaire (section 4).

\courssub{Lire finement : la démarche générale}

\begin{methode}[Technique QQOQCP adaptée à la lecture]
Avant de répondre aux questions, interroge le texte avec les six interrogatifs espagnols :
\begin{enumerate}
\item \esp{¿Quién?} Qui parle ? Qui est concerné ? (l'auteur, \esp{los países del Norte}, \esp{los países del Sur}...)
\item \esp{¿Qué?} De quoi parle-t-on exactement ? (\esp{la cooperación}, \esp{la deuda}, \esp{los intercambios}...)
\item \esp{¿Dónde?} Où se situe l'action ? (\esp{África}, \esp{América Latina}, \esp{las instituciones internacionales}...)
\item \esp{¿Cuándo?} À quelle époque ? (\esp{desde los años setenta}, \esp{hoy en día}...)
\item \esp{¿Cómo?} Comment le phénomène se manifeste-t-il ? (\esp{mediante la ayuda}, \esp{el comercio}, \esp{los acuerdos}...)
\item \esp{¿Por qué?} Pourquoi ? Quelles causes, quelles intentions ? (\esp{para reducir la desigualdad}, \esp{para convencer}...)
\end{enumerate}
Réponds mentalement à ces six questions avant de rédiger : tu auras déjà 80 \% de la compréhension.
\end{methode}

\begin{attention}[Éviter la traduction mot à mot]
Le piège classique du francophone est de traduire chaque mot séparément avant de comprendre la phrase. Résultat : des contresens. Privilégie toujours le \cle{sens global} de la phrase ou du paragraphe. Ainsi \esp{a pesar de los esfuerzos} ne se traduit pas « à pesar des efforts » mais « malgré les efforts » ; \esp{llevar a cabo} ne veut pas dire « porter à la barque » mais « mener à bien, réaliser ». Lis d'abord toute la phrase, identifie le verbe et son sujet, puis seulement après précise le sens des mots difficiles grâce au contexte.
\end{attention}

\courssub{Le texte support (rappel)}

Pour travailler cette section, relisons le texte étudié dans la section 1.

\begin{textebox}[Las relaciones Norte-Sur : entre cooperación y desigualdad]
Desde los años setenta, los informes de la CNUCED denuncian la desigualdad entre los países del Norte y los del Sur. A pesar de los esfuerzos de la cooperación internacional, la brecha entre ricos y pobres sigue siendo profunda. Según los datos disponibles, una parte importante de la población mundial vive con menos de dos dólares al día, mientras que los países desarrollados concentran la mayor parte de la riqueza.

Sin embargo, la cooperación no es una idea muerta. Organizaciones como la ONU o la Unión Africana promueven programas de desarrollo en África y en América Latina. Por ejemplo, en Camerún, varios proyectos financiados por socios internacionales han mejorado el acceso al agua potable en las zonas rurales. Además, la ayuda humanitaria salva vidas cada año durante las crisis alimentarias.

No obstante, algunos economistas critican esta ayuda : según ellos, crea dependencia y no resuelve los problemas estructurales, como la deuda externa o los precios injustos de las materias primas. El cacao camerunés, por ejemplo, se vende barato en el mercado mundial y se transforma en el Norte, donde genera beneficios muy superiores.

Por lo tanto, muchos países del Sur reclaman hoy un comercio más justo en lugar de una simple ayuda. Como dijo un dirigente africano : « No queremos limosna, queremos justicia ». La globalización, en efecto, solo será positiva si respeta los derechos de todos los pueblos y si los ciudadanos, en el Norte como en el Sur, exigen a sus gobiernos políticas más solidarias.
\end{textebox}

\textbf{Glossaire.} \esp{la brecha} : le fossé ; \esp{el acceso al agua potable} : l'accès à l'eau potable ; \esp{las materias primas} : les matières premières ; \esp{la limosna} : l'aumône ; \esp{reclamar} : réclamer ; \esp{solidario} : solidaire.

\courssub{Les neuf questions de compréhension}

Voici la batterie complète de questions qui guident la lecture détaillée : trois questions \cle{globales} (idée principale, thèse, public visé) et six questions \cle{fines} (chiffres, exemples, arguments).

\begin{textebox}[Las nueve preguntas de comprensión]
\textbf{Comprensión global}

1. ¿Cuál es la idea principal del texto ? \emph{(Quelle est l'idée principale du texte ?)}

2. ¿Cuál es la tesis del autor sobre la cooperación internacional ? \emph{(Quelle est la thèse de l'auteur sur la coopération internationale ?)}

3. ¿A qué público se dirige este artículo ? \emph{(À quel public cet article s'adresse-t-il ?)}

\textbf{Comprensión fina}

4. ¿Qué datos numéricos se mencionan sobre la pobreza mundial ? \emph{(Quelles données chiffrées sont mentionnées sur la pauvreté mondiale ?)}

5. ¿Qué países o regiones se citan como ejemplos ? \emph{(Quels pays ou régions sont cités comme exemples ?)}

6. ¿Qué organizaciones internacionales aparecen en el texto y qué papel desempeñan ? \emph{(Quelles organisations internationales apparaissent et quel rôle jouent-elles ?)}

7. ¿Qué argumentos se dan a favor de la cooperación ? \emph{(Quels arguments sont donnés en faveur de la coopération ?)}

8. ¿Qué críticas se hacen a la ayuda internacional ? \emph{(Quelles critiques sont faites de l'aide internationale ?)}

9. ¿Qué solución proponen los países del Sur según el último párrafo ? \emph{(Quelle solution les pays du Sud proposent-ils selon le dernier paragraphe ?)}
\end{textebox}

\begin{popfigure}
\begin{tikzpicture}
  \node[draw,rounded corners,fill=popblueL,align=center] (q1) {Questions\\globales (3)};
  \node[draw,rounded corners,fill=poporangeL,align=center,right=1.5cm of q1] (q2) {Questions\\fines (6)};
  \node[draw,rounded corners,fill=popgreenL,align=center,below=1cm of q2] (r) {Reformulation\\personnelle};
  \draw[->,thick,popdark] (q1) -- (q2);
  \draw[->,thick,popdark] (q2) -- (r);
  \draw[->,thick,popdark] (q1) -- (r);
\end{tikzpicture}
\legende{Progression de la lecture détaillée : du sens global vers la reformulation personnelle.}
\end{popfigure}

\courssub{Correction commentée des questions}

\begin{exoresolu}[Questions globales : modèles de réponses]
Réponds en espagnol, en phrases complètes, aux questions 1 à 3.
\tcblower
\textbf{1. Idea principal :} \esp{El texto trata de la desigualdad entre los países del Norte y del Sur y del papel de la cooperación internacional.} « Le texte traite de l'inégalité entre les pays du Nord et du Sud et du rôle de la coopération internationale. »

\textbf{2. Tesis del autor :} \esp{El autor piensa que la ayuda no basta : hace falta un comercio más justo y políticas solidarias.} « L'auteur pense que l'aide ne suffit pas : il faut un commerce plus juste et des politiques solidaires. »

\textbf{3. Público visado :} \esp{El artículo se dirige a un público general, sobre todo a los jóvenes lectores y a los ciudadanos interesados por los problemas mundiales.} « L'article s'adresse à un public général, surtout aux jeunes lecteurs et aux citoyens intéressés par les problèmes mondiaux. »

\textbf{Méthode :} pour l'idée principale, cherche le thème (de quoi parle-t-on ?) + le point de vue (qu'en dit-on ?). Pour le public, observe le ton (informatif, engagé) et l'absence de vocabulaire trop technique.
\end{exoresolu}

\begin{exoresolu}[Questions fines : modèles de réponses]
Réponds aux questions 4 à 9 en citant le texte.
\tcblower
\textbf{4. Datos numéricos :} \esp{El texto menciona los años setenta y la cifra de menos de dos dólares al día para una parte importante de la población mundial.} « Le texte mentionne les années soixante-dix et le chiffre de moins de deux dollars par jour. »

\textbf{5. Países y regiones :} \esp{Se citan África, América Latina y, como ejemplo concreto, Camerún, con sus proyectos de agua potable y su cacao.} « Sont cités l'Afrique, l'Amérique latine et, comme exemple concret, le Cameroun, avec ses projets d'eau potable et son cacao. »

\textbf{6. Organizaciones :} \esp{La CNUCED denuncia la desigualdad ; la ONU y la Unión Africana promueven programas de desarrollo.} « La CNUCED dénonce l'inégalité ; l'ONU et l'Union africaine promeuvent des programmes de développement. »

\textbf{7. Argumentos a favor :} \esp{La cooperación mejora el acceso al agua potable y la ayuda humanitaria salva vidas durante las crisis alimentarias.} « La coopération améliore l'accès à l'eau potable et l'aide humanitaire sauve des vies pendant les crises alimentaires. »

\textbf{8. Críticas :} \esp{Según algunos economistas, la ayuda crea dependencia y no resuelve los problemas estructurales como la deuda externa y los precios injustos de las materias primas.} « Selon certains économistes, l'aide crée une dépendance et ne résout pas les problèmes structurels comme la dette extérieure et les prix injustes des matières premières. »

\textbf{9. Solución propuesta :} \esp{Los países del Sur reclaman un comercio más justo en lugar de una simple ayuda.} « Les pays du Sud réclament un commerce plus juste au lieu d'une simple aide. »
\end{exoresolu}

\begin{attention}[Citer le texte sans recopier]
Dans une réponse de compréhension, il faut \emph{s'appuyer} sur le texte mais le \emph{reformuler}. Évite de recopier trois lignes entières : extrais l'information et intègre-la dans ta propre phrase. Utilise des verbes comme \esp{el autor afirma que...}, \esp{según el texto...}, \esp{el artículo denuncia...}. Et attention aux faux amis : \esp{actualmente} signifie « actuellement, en ce moment » (et non « actuellement » au sens de « en fait », qui se dit \esp{en realidad}) ; \esp{una cifra} signifie « un chiffre, un montant » (et non « un chiffre d'affaires », qui se dit \esp{el volumen de negocios}).
\end{attention}

\courssub{Exercice de reformulation : résumer chaque paragraphe en une phrase}

Résumer un paragraphe en une seule phrase est un excellent entraînement à la synthèse, compétence exigée au commentaire de texte. La règle : \textbf{une phrase = sujet + verbe + idée essentielle}, sans exemple ni chiffre secondaire.

\begin{exemplebox}[Modèles de résumés par paragraphe]
\begin{itemize}
\item \textbf{Paragraphe 1 :} \esp{A pesar de la cooperación, la desigualdad entre el Norte y el Sur sigue siendo muy profunda.} « Malgré la coopération, l'inégalité entre le Nord et le Sud reste très profonde. »
\item \textbf{Paragraphe 2 :} \esp{La cooperación internacional consigue resultados concretos, como el acceso al agua potable en Camerún.} « La coopération internationale obtient des résultats concrets, comme l'accès à l'eau potable au Cameroun. »
\item \textbf{Paragraphe 3 :} \esp{Algunos economistas critican la ayuda porque crea dependencia y no ataca los problemas estructurales.} « Certains économistes critiquent l'aide parce qu'elle crée une dépendance et ne s'attaque pas aux problèmes structurels. »
\item \textbf{Paragraphe 4 :} \esp{Los países del Sur prefieren un comercio justo a la limosna y exigen políticas solidarias.} « Les pays du Sud préfèrent un commerce juste à l'aumône et exigent des politiques solidaires. »
\end{itemize}
\end{exemplebox}

\begin{exercice}[À ton tour]
Résume en une phrase espagnole le sens global de tout le texte (maximum 25 mots). Commence par : \esp{El texto muestra que...}
\end{exercice}

\begin{corrige}[Exercice : résumé global]
Proposition de corrigé : \esp{El texto muestra que la cooperación Norte-Sur, aunque útil, no basta para acabar con la desigualdad y que los países del Sur reclaman un comercio más justo.} « Le texte montre que la coopération Nord-Sud, bien qu'utile, ne suffit pas à éliminer l'inégalité et que les pays du Sud réclament un commerce plus juste. » Toute formulation correcte qui contient les deux idées (utilité de l'aide + insuffisance / exigence de justice commerciale) est acceptable.
\end{corrige}

\courssub{Les connecteurs logiques du texte et leur fonction}

Repérer les connecteurs, c'est découvrir la \cle{charpente argumentative} du texte. Voici ceux du texte support, classés par fonction :

\begin{center}
\begin{tabularx}{\linewidth}{|X|X|X|}
\hline
\rowcolor{popblueL} Conector & Función & Traduction \\
\hline
\esp{a pesar de} & concesión & malgré \\
\hline
\esp{sin embargo} & oposición & cependant \\
\hline
\esp{mientras que} & contraste & tandis que \\
\hline
\esp{no obstante} & oposición & néanmoins \\
\hline
\esp{por ejemplo} & ejemplificación & par exemple \\
\hline
\esp{además} & adición & de plus \\
\hline
\esp{según} & fuente / opinión & selon \\
\hline
\esp{por lo tanto} & consecuencia & par conséquent \\
\hline
\esp{en lugar de} & sustitución & au lieu de \\
\hline
\esp{en efecto} & confirmación & en effet \\
\hline
\esp{como} (causal) & causa & comme, parce que \\
\hline
\end{tabularx}
\end{center}

\begin{propriete}[La structure argumentative du texte]
Le texte suit un plan dialectique classique en trois mouvements, signalés par les connecteurs :
\begin{enumerate}
\item \textbf{Thèse implicite du problème} : \esp{la desigualdad persiste} (\esp{a pesar de}).
\item \textbf{Arguments en faveur de la cooperación} : \esp{sin embargo} ouvre le paragraphe 2 (la coopération n'est pas une idée morte).
\item \textbf{Antithèse et synthèse} : \esp{no obstante} introduit les critiques (paragraphe 3), puis \esp{por lo tanto} annonce la conclusion : un commerce plus juste (paragraphe 4).
\end{enumerate}
Au commentaire, citer ces connecteurs prouve que tu as compris la progression de l'argumentation.
\end{propriete}

\begin{exemplebox}[Exemples traduits : les connecteurs en contexte]
\begin{itemize}
\item \esp{A pesar de los esfuerzos, la brecha sigue siendo profunda.} « Malgré les efforts, le fossé reste profond. »
\item \esp{Sin embargo, la cooperación no es una idea muerta.} « Cependant, la coopération n'est pas une idée morte. »
\item \esp{Los países ricos concentran la riqueza, mientras que el Sur vive en la pobreza.} « Les pays riches concentrent la richesse, tandis que le Sud vit dans la pauvreté. »
\item \esp{No obstante, algunos economistas critican esta ayuda.} « Néanmoins, certains économistes critiquent cette aide. »
\item \esp{Por lo tanto, reclaman un comercio más justo.} « Par conséquent, ils réclament un commerce plus juste. »
\end{itemize}
\end{exemplebox}

\begin{attention}[Pièges des connecteurs pour les francophones]
\begin{itemize}
\item \esp{Sin embargo} se place en début de phrase, suivi d'une virgule : \esp{Sin embargo, ...} Ne le confonds pas avec \esp{sin que} (qui exige le subjonctif : \esp{sin que lo sepan}, « sans qu'ils le sachent »).
\item \esp{A pesar de} est suivi d'un \textbf{nom} ou d'un \textbf{infinitif} : \esp{a pesar de los esfuerzos}, \esp{a pesar de ser pobre}. Jamais d'une phrase conjuguée directement : il faut alors \esp{a pesar de que} + indicatif ou subjonctif (\esp{a pesar de que llueve}, « bien qu'il pleuve »).
\item Ne traduis pas \esp{por lo tanto} par « pour le tant » : c'est « par conséquent ». Et \esp{en efecto} signifie bien « en effet, effectivement », sans nuance ironique.
\end{itemize}
\end{attention}

\courssub{Informations implicites : ton et intention de l'auteur}

Un bon lecteur lit aussi \emph{entre les lignes}. Deux indices permettent de dégager l'implicite :

\begin{itemize}
\item \textbf{Le ton.} Le mélange de données objectives (\esp{menos de dos dólares al día}) et de formules engagées (\esp{precios injustos}, \esp{no queremos limosna, queremos justicia}) révèle un ton à la fois \cle{informatif} et \cle{engagé}, voire militant. L'auteur ne se contente pas de décrire : il dénonce (\esp{denuncian}, \esp{reclaman}, \esp{exigen}).
\item \textbf{L'intention persuasive.} La citation finale du dirigeant africain et l'appel aux citoyens (\esp{los ciudadanos... exigen a sus gobiernos}) montrent que l'auteur cherche à \cle{convaincre} le lecteur et à le pousser à l'action citoyenne, pas seulement à l'informer. C'est la fonction \emph{conative} du langage.
\end{itemize}

\begin{exemplebox}[Questions sur l'implicite et réponses attendues]
\esp{¿Cuál es el tono del texto ?} → \esp{El tono es informativo pero también comprometido : el autor denuncia la injusticia.} « Le ton est informatif mais aussi engagé : l'auteur dénonce l'injustice. »

\esp{¿Qué intención tiene el autor ?} → \esp{Su intención es persuadir al lector de que la ayuda no basta y de que hay que exigir un comercio justo.} « Son intention est de persuader le lecteur que l'aide ne suffit pas et qu'il faut exiger un commerce juste. »
\end{exemplebox}

\begin{savaistu}[D'où vient l'expression « Norte-Sur » ?]
Le terme \esp{Norte-Sur} apparaît dans les rapports de la \esp{CNUCED} (Conférence des Nations unies sur le commerce et le développement) depuis les années 1970. Il remplace peu à peu l'expression « Tiers-Monde » (\esp{Tercer Mundo}), jugée trop péjorative. En français, on parle du « rapport Brandt » (1980), qui a popularisé la ligne imaginaire séparant les pays riches du « Nord » des pays pauvres du « Sud » — même si l'Australie, au sud géographique, appartient au « Nord » économique ! En espagnol, on dit aussi \esp{países en vías de desarrollo} (pays en voie de développement).
\end{savaistu}

\courssub{Auto-évaluation : grille de critères}

Après avoir répondu aux neuf questions et rédigé tes résumés, évalue-toi honnêtement avec la grille suivante. Note chaque critère de 0 à 2 : 0 = non acquis, 1 = en cours, 2 = acquis.

\begin{center}
\begin{tabularx}{\linewidth}{|X|X|X|}
\hline
\rowcolor{popblueL} Critère & Indicateur de réussite & Note (0-2) \\
\hline
\cle{Justesse} & Mes réponses correspondent aux informations du texte, sans contresens. & \\
\hline
\cle{Vocabulaire} & J'emploie le lexique du thème : \esp{cooperación}, \esp{deuda}, \esp{desarrollo}, \esp{brecha}. & \\
\hline
\cle{Cohérence} & Mes phrases sont liées par des connecteurs corrects (\esp{sin embargo}, \esp{por lo tanto}). & \\
\hline
\cle{Reformulación} & Je résume sans recopier le texte mot à mot. & \\
\hline
\cle{Gramática} & Concordance des temps, accords, accents respectés. & \\
\hline
\cle{Comprensión implícita} & J'ai identifié le ton et l'intention persuasive de l'auteur. & \\
\hline
\end{tabularx}
\end{center}

\begin{aretenir}[L'essentiel de la compréhension détaillée]
\begin{itemize}
\item Interroge le texte avec la technique \esp{QQOQCP} avant de répondre.
\item Distingue toujours questions \cle{globales} (idée, thèse, public) et questions \cle{fines} (chiffres, exemples, arguments).
\item Reformule au lieu de recopier ; cite le texte avec \esp{según el autor...}.
\item Les connecteurs (\esp{a pesar de}, \esp{sin embargo}, \esp{no obstante}, \esp{por lo tanto}) révèlent la structure argumentative.
\item Cherche l'implicite : ton engagé, intention persuasive, appel au lecteur.
\end{itemize}
\end{aretenir}

Cette lecture fine du texte nous a permis d'en extraire toute la substance : informations explicites, arguments, connecteurs et intentions cachées. Nous sommes maintenant prêts pour la section suivante, qui organisera le \cle{lexique thématique} des relations internationales et des systèmes politiques en champs sémantiques clairs, avant de passer à la méthode du commentaire de texte proprement dit.

\coursec{Lexique thématique par champs}

Dans la section précédente, nous avons lu et analysé en détail le texte support sur les relations Nord-Sud. Nous avons rencontré des mots comme \esp{cooperación}, \esp{deuda}, \esp{democracia} ou \esp{derechos humanos}. Pour commenter un texte politique à l'examen, il ne suffit pas de reconnaître ces mots : il faut les \cle{maîtriser activement}, c'est-à-dire connaître leur genre, leur catégorie grammaticale, leurs collocations (les mots avec lesquels ils s'emploient) et savoir les réutiliser dans une phrase correcte. C'est l'objectif de cette section : organiser le vocabulaire en quatre champs lexicaux clairs, le fixer par des exemples traduits et des exercices.

\begin{popfigure}
\begin{tikzpicture}
  \node[draw,rounded corners,fill=popblueL] (c1) {Organisations};
  \node[draw,rounded corners,fill=poporangeL,right=1.5cm of c1] (c2) {Économie};
  \node[draw,rounded corners,fill=popgreenL,below=1cm of c1] (c3) {Systèmes};
  \node[draw,rounded corners,fill=poppinkL,right=1.5cm of c3] (c4) {Droits};
  \draw[->,thick,popblue] (c1) -- (c2);
  \draw[->,thick,popgreen] (c3) -- (c4);
  \draw[->,thick,popmuted] (c1) -- (c3);
  \draw[->,thick,popmuted] (c2) -- (c4);
\end{tikzpicture}
\legende{Les quatre champs lexicaux de la vie internationale : les organisations agissent sur l'économie ; les systèmes politiques garantissent (ou non) les droits.}
\end{popfigure}

\courssub{Champ 1 : Las organizaciones internacionales}

Observons d'abord ces phrases, typiques de la presse hispanophone :

\esp{La ONU ha pedido un alto el fuego inmediato.} « L'ONU a demandé un cessez-le-feu immédiat. »

\esp{La Unión Africana promueve la integración del continente.} « L'Union africaine promeut l'intégration du continent. »

\esp{El FMI concede préstamos a los países endeudados.} « Le FMI accorde des prêts aux pays endettés. »

\begin{definition}[Las organizaciones internacionales]
Une \cle{organisation internationale} est une institution créée par un accord entre plusieurs États pour coopérer dans un domaine (paix, économie, santé). En espagnol, leurs sigles (\esp{siglas}) sont généralement féminins quand le mot sous-entendu est \esp{organización} : \esp{la ONU}, \esp{la UA}, \esp{la UE}.
\end{definition}

\begin{center}
\begin{tabularx}{\linewidth}{|l|l|X|}
\hline
\rowcolor{popblueL} Español & Français & Ejemplo \\
\hline
la ONU (Organización de las Naciones Unidas), n.f. & l'ONU & \esp{La ONU fue fundada en 1945.} (L'ONU fut fondée en 1945.) \\
\hline
la UA (Unión Africana), n.f. & l'UA (Union africaine) & \esp{La UA tiene su sede en Addis Abeba.} (L'UA a son siège à Addis-Abeba.) \\
\hline
la UE (Unión Europea), n.f. & l'UE (Union européenne) & \esp{España forma parte de la UE desde 1986.} (L'Espagne fait partie de l'UE depuis 1986.) \\
\hline
el FMI (Fondo Monetario Internacional), n.m. & le FMI & \esp{El FMI exige reformas económicas.} (Le FMI exige des réformes économiques.) \\
\hline
el BM (Banco Mundial), n.m. & la Banque mondiale & \esp{El BM financia proyectos de desarrollo.} (La Banque mondiale finance des projets de développement.) \\
\hline
la cooperación, n.f. & la coopération & \esp{La cooperación entre Camerún y España es estrecha.} (La coopération entre le Cameroun et l'Espagne est étroite.) \\
\hline
el tratado, n.m. & le traité & \esp{Los países firmaron un tratado de paz.} (Les pays ont signé un traité de paix.) \\
\hline
el/la embajador(a), n. & l'ambassadeur / l'ambassadrice & \esp{La embajadora presentó sus cartas credenciales.} (L'ambassadrice a présenté ses lettres de créance.) \\
\hline
\end{tabularx}
\end{center}

\begin{propriete}[Sigles et article]
Les sigles d'organisations s'écrivent en \cle{majuscules} : \esp{ONU}, \esp{UA}, \esp{UE}, \esp{FMI}. Deux usages coexistent : avec article quand on prononce le sigle comme un nom (\esp{la ONU}, \esp{el FMI}), sans article dans les titres de presse ou les formules administratives (\esp{ONU pide calma}). À l'écrit scolaire, utilisez toujours l'article : \esp{la ONU}, \esp{la UA}. Attention au genre : \esp{la ONU} (car \esp{organización}) mais \esp{el FMI} (car \esp{fondo}) et \esp{el BM} (car \esp{banco}).
\end{propriete}

\begin{propriete}[El Norte y el Sur]
Dans les textes sur les relations internationales, \esp{el Norte} désigne les pays développés et \esp{el Sur} les pays en voie de développement : ce sont des noms communs employés avec majuscule et article dans ce sens géopolitique. On dit : \esp{las relaciones Norte-Sur} (les relations Nord-Sud), \esp{la brecha entre el Norte y el Sur} (le fossé entre le Nord et le Sud), \esp{los países del Sur} (les pays du Sud).
\end{propriete}

\begin{savaistu}[La Guinée équatoriale, pays hispanophone]
La Guinée équatoriale, voisine du Cameroun, est le seul pays d'Afrique subsaharienne où l'espagnol est langue officielle. Elle est membre de l'Union africaine et des Nations unies : un excellent exemple à citer dans vos copies pour montrer que l'espagnol est aussi une langue africaine.
\end{savaistu}

\courssub{Champ 2 : Los conceptos económicos}

\begin{definition}[Conceptos económicos clave]
\esp{Desarrollo} (n.m., développement) : progrès économique et social d'un pays. \esp{Deuda} (n.f., dette) : somme d'argent qu'un pays doit rembourser. \esp{Globalización} (n.f., mondialisation) : interconnexion croissante des économies et des cultures du monde.
\end{definition}

\begin{center}
\begin{tabularx}{\linewidth}{|l|l|X|}
\hline
\rowcolor{poporangeL} Español & Français & Ejemplo \\
\hline
el desarrollo, n.m. & le développement & \esp{El desarrollo sostenible es un objetivo mundial.} (Le développement durable est un objectif mondial.) \\
\hline
la deuda (externa), n.f. & la dette (extérieure) & \esp{Muchos países piden la cancelación de la deuda.} (Beaucoup de pays demandent l'annulation de la dette.) \\
\hline
la globalización, n.f. & la mondialisation & \esp{La globalización une y divide al mismo tiempo.} (La mondialisation unit et divise à la fois.) \\
\hline
el comercio (internacional), n.m. & le commerce (international) & \esp{El comercio de cacao es vital para Camerún.} (Le commerce du cacao est vital pour le Cameroun.) \\
\hline
la inversión (extranjera), n.f. & l'investissement (étranger) & \esp{La inversión extranjera crea empleo.} (L'investissement étranger crée de l'emploi.) \\
\hline
la ayuda al desarrollo, n.f. & l'aide au développement & \esp{La ayuda no siempre llega a los más pobres.} (L'aide n'arrive pas toujours aux plus pauvres.) \\
\hline
el intercambio, n.m. & l'échange & \esp{El intercambio comercial beneficia a ambas partes.} (L'échange commercial profite aux deux parties.) \\
\hline
la pobreza / la riqueza, n.f. & la pauvreté / la richesse & \esp{La lucha contra la pobreza es prioritaria.} (La lutte contre la pauvreté est prioritaire.) \\
\hline
\end{tabularx}
\end{center}

\begin{attention}[Genres trompeurs]
\esp{Deuda} est féminin : \esp{la deuda}, \esp{una deuda enorme}. Mais \esp{déficit} est masculin : \esp{el déficit}. De même, \esp{el desarrollo} est masculin : ne vous fiez pas au français, vérifiez toujours la terminaison et l'article. Retenez le trio : \esp{la deuda}, \esp{el déficit}, \esp{el desarrollo}. Autre précision : \esp{la inversión} signifie « l'investissement » au sens financier ; c'est le seul sens utile dans un texte économique, le contexte ne prête donc pas à confusion.
\end{attention}

\begin{savaistu}[Un mot venu de l'anglais]
Le mot \esp{globalización} vient de l'anglais \emph{globalization}, popularisé dans les années 1980. L'espagnol a adapté la graphie (\esp{-ización} avec accent sur la \esp{ó}) et le genre (féminin, comme tous les mots en \esp{-ción}). On trouve aussi le verbe \esp{globalizar} et l'adjectif \esp{globalizado/a} : \esp{una economía globalizada}, « une économie mondialisée ».
\end{savaistu}

\courssub{Champ 3 : Los sistemas políticos}

\begin{definition}[Democracia]
\esp{Democracia} (n.f.) : sistema político en el que el poder reside en el pueblo, que lo ejerce directamente o por medio de representantes elegidos. « Démocratie : système politique dans lequel le pouvoir réside dans le peuple, qui l'exerce directement ou par l'intermédiaire de représentants élus. »
\end{definition}

\begin{center}
\begin{tabularx}{\linewidth}{|l|l|X|}
\hline
\rowcolor{popgreenL} Español & Français & Ejemplo \\
\hline
la democracia, n.f. & la démocratie & \esp{En una democracia, el pueblo elige a sus gobernantes.} (Dans une démocratie, le peuple élit ses gouvernants.) \\
\hline
la dictadura, n.f. & la dictature & \esp{La dictadura suprime las libertades.} (La dictature supprime les libertés.) \\
\hline
la monarquía (parlamentaria), n.f. & la monarchie (parlementaire) & \esp{España es una monarquía parlamentaria.} (L'Espagne est une monarchie parlementaire.) \\
\hline
la república, n.f. & la république & \esp{Camerún es una república.} (Le Cameroun est une république.) \\
\hline
el federalismo, n.m. & le fédéralisme & \esp{Alemania tiene un sistema federal.} (L'Allemagne a un système fédéral.) \\
\hline
el gobierno, n.m. & le gouvernement & \esp{El gobierno presentó su programa.} (Le gouvernement a présenté son programme.) \\
\hline
el parlamento, n.m. & le parlement & \esp{El parlamento vota las leyes.} (Le parlement vote les lois.) \\
\hline
las elecciones, n.f.pl. & les élections & \esp{Las elecciones serán en octubre.} (Les élections auront lieu en octobre.) \\
\hline
el partido político, n.m. & le parti politique & \esp{Hay varios partidos políticos en el país.} (Il y a plusieurs partis politiques dans le pays.) \\
\hline
el/la ciudadano/a, n. & le/la citoyen(ne) & \esp{Los ciudadanos votan cada cinco años.} (Les citoyens votent tous les cinq ans.) \\
\hline
\end{tabularx}
\end{center}

\begin{attention}[Faux amis et calques]
\begin{itemize}
\item \esp{El gobierno} = le gouvernement ; ne dites jamais \esp{governación}, qui n'existe pas.
\item \esp{Las elecciones} s'emploie presque toujours au pluriel en espagnol ; « une élection » partielle = \esp{una elección}.
\item \esp{El partido} = le parti politique ; ne pas confondre avec \esp{el partido} de football (le match) : le contexte décide.
\item « Le pouvoir » se traduit \esp{el poder} (sans u) : \esp{el poder reside en el pueblo}.
\item « Élire » se dit \esp{elegir} (verbe irrégulier : \esp{elijo, eliges, elige\ldots}) ; « les élus » = \esp{los elegidos} ou \esp{los representantes elegidos}.
\end{itemize}
\end{attention}

\courssub{Champ 4 : Derechos, libertades y deberes}

\begin{definition}[Derechos humanos]
\esp{Los derechos humanos} (n.m.pl.) : ensemble des droits fondamentaux de tout être humain (vie, liberté, égalité, éducation), proclamés par la Déclaration universelle de 1948 : \esp{la Declaración Universal de los Derechos Humanos}.
\end{definition}

\begin{center}
\begin{tabularx}{\linewidth}{|l|l|X|}
\hline
\rowcolor{poppinkL} Español & Français & Ejemplo \\
\hline
los derechos humanos, n.m.pl. & les droits de l'homme & \esp{Los derechos humanos son universales.} (Les droits de l'homme sont universels.) \\
\hline
el sufragio (universal), n.m. & le suffrage (universel) & \esp{El sufragio universal es un derecho conquistado.} (Le suffrage universel est un droit conquis.) \\
\hline
la libertad de expresión, n.f. & la liberté d'expression & \esp{La libertad de expresión permite criticar al gobierno.} (La liberté d'expression permet de critiquer le gouvernement.) \\
\hline
la igualdad, n.f. & l'égalité & \esp{La igualdad entre hombres y mujeres es esencial.} (L'égalité entre hommes et femmes est essentielle.) \\
\hline
la justicia, n.f. & la justice & \esp{La justicia debe ser independiente.} (La justice doit être indépendante.) \\
\hline
el deber, n.m. & le devoir & \esp{Pagar los impuestos es un deber ciudadano.} (Payer les impôts est un devoir citoyen.) \\
\hline
el voto / votar, n.m. / v. & le vote / voter & \esp{El voto es secreto y libre.} (Le vote est secret et libre.) \\
\hline
la constitución, n.f. & la constitution & \esp{La constitución garantiza los derechos.} (La constitution garantit les droits.) \\
\hline
\end{tabularx}
\end{center}

\begin{propriete}[Derecho : un mot, plusieurs sens]
\esp{El derecho} (masculin) signifie « le droit » au sens juridique : \esp{el derecho a la educación} (le droit à l'éducation), \esp{el derecho internacional} (le droit international). Ne pas confondre avec \esp{derecho/a} adjectif (« droit », opposé de gauche : \esp{la mano derecha}, la main droite) ni avec l'adverbe \esp{derecho} (« tout droit » : \esp{siga derecho}, continuez tout droit). Dans notre champ lexical, c'est toujours le nom masculin juridique. Notez la construction : \esp{el derecho a} + nom ou infinitif : \esp{el derecho a votar} (le droit de voter), \esp{el derecho a la vida} (le droit à la vie).
\end{propriete}

\begin{exemplebox}[Exemples traduits à réutiliser]
\esp{La Unión Africana promueve la integración.} → « L'Union africaine promeut l'intégration. »

\esp{Todo ciudadano tiene derechos, pero también deberes.} → « Tout citoyen a des droits, mais aussi des devoirs. »

\esp{En una democracia, la libertad de prensa está protegida por la ley.} → « Dans une démocratie, la liberté de la presse est protégée par la loi. »

\esp{La deuda externa frena el desarrollo de muchos países africanos.} → « La dette extérieure freine le développement de beaucoup de pays africains. »
\end{exemplebox}

\courssub{Texte d'application : un discours citoyen}

Lisons maintenant un court texte original qui réutilise le lexique des quatre champs. Repérez chaque mot étudié.

\begin{textebox}[Un joven diputado habla a los electores]
Queridos conciudadanos: vivimos en un mundo globalizado, donde las decisiones de la ONU o del FMI influyen en nuestra vida diaria, desde el precio del cacao hasta el empleo de los jóvenes. Nuestro país, miembro de la Unión Africana, cree en la cooperación entre los pueblos, pero también exige relaciones más justas entre el Norte y el Sur. La deuda externa no debe impedir el desarrollo de nuestras escuelas y de nuestros hospitales.

Sin embargo, la solución no vendrá solamente del exterior. Depende de nosotros, los ciudadanos. En una democracia, el poder reside en el pueblo: el sufragio universal nos permite elegir a nuestros representantes, y la libertad de expresión nos permite criticarlos cuando se equivocan. La igualdad ante la ley y una justicia independiente son el fundamento de la república.

Pero los derechos van acompañados de deberes. Votar es un derecho; informarse antes de votar es un deber. Expresarse es un derecho; respetar la opinión del otro es un deber. Pagar los impuestos, proteger el medio ambiente, educar a nuestros hijos: todo eso construye la nación.

Les pido, pues, que participen en la vida pública. Un pueblo que vota, que debate y que respeta la ley es un pueblo libre. La globalización puede ser una oportunidad si los ciudadanos están despiertos. ¡Gracias, y viva la democracia!
\end{textebox}

\textbf{Glossaire :} \esp{conciudadanos} = concitoyens ; \esp{el empleo} = l'emploi ; \esp{exigir} = exiger ; \esp{impedir} = empêcher ; \esp{equivocarse} = se tromper ; \esp{el fundamento} = le fondement ; \esp{los impuestos} = les impôts ; \esp{el medio ambiente} = l'environnement ; \esp{despierto} = éveillé, vigilant.

\textbf{Questions de repérage lexical :}
\begin{enumerate}
\item Relevez dans le texte trois noms ou sigles d'organisations internationales.
\item Trouvez deux noms de systèmes politiques et trois droits mentionnés.
\item Quels \emph{deberes} (devoirs) l'orateur cite-t-il ? Donnez-en trois.
\end{enumerate}

\begin{corrige}[Questions de repérage lexical]
1. Organisations : \esp{la ONU}, \esp{el FMI}, \esp{la Unión Africana}. 2. Systèmes politiques : \esp{la democracia}, \esp{la república}. Droits : \esp{el sufragio universal}, \esp{la libertad de expresión}, \esp{la igualdad ante la ley} (on peut ajouter \esp{votar} et \esp{expresarse}). 3. Devoirs : \esp{informarse antes de votar}, \esp{respetar la opinión del otro}, \esp{pagar los impuestos}, \esp{proteger el medio ambiente}, \esp{educar a nuestros hijos} (trois suffisent).
\end{corrige}

\courssub{Exercices de fixation}

\begin{exoresolu}[Correspondance mot / définition]
\textbf{Énoncé.} Associez chaque mot espagnol à sa définition : 1. \esp{deuda} ; 2. \esp{sufragio} ; 3. \esp{dictadura} ; 4. \esp{inversión} ; 5. \esp{tratado}. Définitions : a. sistema sin libertades donde manda una sola persona o grupo ; b. dinero que un país debe devolver ; c. acuerdo firmado entre Estados ; d. derecho a votar ; e. dinero colocado en un país para producir riqueza.
\tcblower
\textbf{Solution détaillée.} 1-b : \esp{la deuda} est l'argent qu'un pays doit rembourser. 2-d : \esp{el sufragio} est le droit de vote (\esp{sufragio universal}). 3-a : \esp{la dictadura} est le système sans libertés, opposé de \esp{democracia}. 4-e : \esp{la inversión} est l'argent investi pour créer de la richesse. 5-c : \esp{el tratado} est l'accord signé entre États, comme un traité de paix ou de commerce.
\end{exoresolu}

\begin{exercice}[Complétion]
Complétez les phrases avec le mot qui convient (attention au genre et à l'article) : \esp{deuda / democracia / ONU / libertad de expresión / inversión / derechos humanos}.

1. \esp{En una \ldots\ldots, los ciudadanos eligen a sus gobernantes.}

2. \esp{La \ldots\ldots externa de algunos países africanos es muy alta.}

3. \esp{La \ldots\ldots organiza misiones de paz en el mundo.}

4. \esp{Los periodistas defienden la \ldots\ldots.}

5. \esp{La \ldots\ldots extranjera puede crear empleo en Douala.}

6. \esp{La Declaración Universal de los \ldots\ldots fue proclamada en 1948.}
\end{exercice}

\begin{corrige}[Exercice de complétion]
1. \esp{democracia} (féminin : \esp{una democracia}). 2. \esp{deuda} (\esp{la deuda externa}, féminin). 3. \esp{ONU} (\esp{la ONU}, sigle féminin). 4. \esp{libertad de expresión}. 5. \esp{inversión} (\esp{la inversión extranjera}). 6. \esp{derechos humanos} (pluriel masculin).
\end{corrige}

\begin{exercice}[Thème : phrases à traduire en espagnol]
Traduisez en espagnol, en réutilisant le lexique de la leçon :

1. La coopération entre le Cameroun et l'Espagne est très étroite.

2. Dans une démocratie, les citoyens ont des droits et des devoirs.

3. La dette extérieure empêche le développement de nombreux pays du Sud.

4. La liberté d'expression est protégée par la constitution.
\end{exercice}

\begin{corrige}[Exercice de thème]
1. \esp{La cooperación entre Camerún y España es muy estrecha.} (Attention : \esp{Camerún} prend un accent ; \esp{cooperación} est féminin.)

2. \esp{En una democracia, los ciudadanos tienen derechos y deberes.} (\esp{tener} : \esp{tienen}, 3\up{e} personne du pluriel.)

3. \esp{La deuda externa impide el desarrollo de muchos países del Sur.} (\esp{impedir} : \esp{impide}, verbe à changement radical e→i ; \esp{del Sur} avec majuscule.)

4. \esp{La libertad de expresión está protegida por la constitución.} (Passif avec \esp{estar} + participe : état résultant ; accord : \esp{protegida}, féminin singulier.)
\end{corrige}

\begin{attention}[Bilan des pièges à éviter dans vos copies]
\begin{itemize}
\item Genres : \esp{la deuda} mais \esp{el déficit} ; \esp{la ONU} mais \esp{el FMI} ; \esp{el desarrollo}, \esp{el comercio}, \esp{el gobierno} sont masculins.
\item Ne traduisez pas « relations internationales » par une forme incorrecte : le pluriel de \esp{internacional} est \esp{internacionales} (sans accent au pluriel) : \esp{las relaciones internacionales}.
\item \esp{Los derechos humanos} : toujours au pluriel dans cette expression ; « les droits de l'homme » ne se dit jamais \esp{derechos del hombre}.
\item Accentuation : \esp{globalización}, \esp{cooperación}, \esp{inversión}, \esp{expresión}, \esp{elecciones} : tous les mots en \esp{-ción} portent un accent sur la \esp{ó}.
\item Noms de pays : \esp{Camerún} (avec accent), \esp{Guinea Ecuatorial}, \esp{España} (avec \esp{ñ}).
\end{itemize}
\end{attention}

\begin{aretenir}[L'essentiel du lexique]
Quatre champs à connaître : \cle{organisations} (\esp{la ONU, la UA, la UE, el FMI, el BM}), \cle{économie} (\esp{desarrollo, deuda, globalización, comercio, inversión}), \cle{systèmes politiques} (\esp{democracia, dictadura, monarquía, república, federalismo}), \cle{droits et devoirs} (\esp{derechos humanos, sufragio, libertad de expresión, igualdad, justicia, deberes}). Pour chaque mot, retenez le genre, l'article et une phrase modèle : c'est ce triplet qui vous permettra de réussir le commentaire de texte et la prise de position citoyenne de la prochaine section.
\end{aretenir}

\coursec{Méthode du commentaire \& commentaire modèle}

Dans la section précédente, nous avons constitué ensemble un lexique riche et organisé par champs : relations internationales, systèmes politiques, droits et devoirs du citoyen. Ce vocabulaire n'est pas une fin en soi : c'est l'\emph{outil} qui va vous permettre de réussir l'épreuve reine du Bac d'espagnol, le \cle{comentario de texto}. Un commentaire sans lexique précis est vide ; un lexique sans méthode est inutile. Voyons maintenant comment transformer vos connaissances en une copie structurée, argumentée et bien notée.

\courssub{Qu'est-ce qu'un comentario de texto ?}

\begin{definition}[Le commentaire de texte]
Le \cle{comentario de texto} est un exercice d'expression écrite qui consiste à \textbf{analyser} un texte (et non à le résumer) : identifier son thème, dégager la thèse de l'auteur, relever les arguments et les procédés, puis donner une \textbf{opinion personnelle justifiée}. En espagnol, on parle de \esp{analizar, interpretar y valorar un texto}.
\end{definition}

La différence fondamentale avec le résumé (\esp{resumen}) est la suivante : le résumé dit \emph{ce que dit le texte} ; le commentaire explique \emph{comment et pourquoi il le dit}, et \emph{ce que vous en pensez}.

\begin{attention}[Le piège n° 1 : résumer au lieu de commenter]
Ne pas résumer le texte ; il faut \emph{commenter} (analyser, évaluer). Une copie qui paraphrase le texte ligne par ligne (\esp{el autor dice que... luego dice que...}) ne dépasse jamais la moyenne. Chaque paragraphe de votre développement doit contenir : une \textbf{idée analytique}, un \textbf{exemple tiré du texte} (une courte citation entre guillemets) et votre \textbf{appréciation personnelle}.
\end{attention}

\courssub{La structure type du commentaire}

Tout commentaire de texte au Bac suit une architecture en trois mouvements, exactement comme en français ou en philosophie :

\begin{popfigure}
\begin{tikzpicture}[node distance=1.6cm]
  \node[draw=popblue,fill=popblueL,rounded corners,align=center,text width=3.2cm] (int) {\textbf{Introduction}\\accroche, sujet,\\thèse, plan};
  \node[draw=poporange,fill=poporangeL,rounded corners,align=center,text width=3.6cm,right=of int] (dev) {\textbf{Développement}\\2 ou 3 parties :\\argument + exemple};
  \node[draw=popgreen,fill=popgreenL,rounded corners,align=center,text width=3.2cm,right=of dev] (concl) {\textbf{Conclusion}\\bilan,\\ouverture};
  \draw[-{Stealth},thick,popdark] (int) -- (dev);
  \draw[-{Stealth},thick,popdark] (dev) -- (concl);
\end{tikzpicture}
\legende{La structure en trois mouvements du comentario de texto}
\end{popfigure}

\begin{propriete}[Contenu de chaque partie]
\begin{itemize}
\item \textbf{Introduction} (\esp{introducción}) : 4 à 6 lignes. Elle contient une \cle{accroche} (\esp{frase de apertura}), la présentation du \cle{sujet} (\esp{tema}), la \cle{thèse} de l'auteur (\esp{tesis}) et l'annonce du \cle{plan} (\esp{esquema}).
\item \textbf{Développement} (\esp{desarrollo}) : 2 ou 3 parties nettement séparées. Chaque partie = un axe d'analyse, avec arguments, exemples cités du texte et opinion personnelle.
\item \textbf{Conclusion} (\esp{conclusión}) : 3 à 5 lignes. Bilan de l'analyse (\esp{en síntesis}) et ouverture vers une question plus large (\esp{a modo de conclusión, cabe preguntarse...}).
\end{itemize}
\end{propriete}

\begin{methode}[Les 6 étapes du commentaire réussi]
\begin{enumerate}
\item \textbf{Lire \& annoter} : lisez le texte deux fois ; soulignez la thèse, entourez les connecteurs, notez en marge les arguments.
\item \textbf{Formuler la thèse} : rédigez en une phrase ce que l'auteur veut démontrer (\esp{El autor defiende la idea de que...}).
\item \textbf{Choisir 2 ou 3 axes} : par exemple, pour notre texte sur les relations Nord-Sud : la coopération, la conditionnalité de l'aide, la souveraineté des États.
\item \textbf{Chercher des exemples} : pour chaque axe, sélectionnez une citation courte du texte et un exemple du monde réel (ONU, Union africaine, Cameroun, Guinée équatoriale).
\item \textbf{Rédiger} : introduction, développement, conclusion, en employant les connecteurs étudiés.
\item \textbf{Relire} : vérifiez les accords, les accents, les subjonctifs et la ponctuation espagnole (¿...? ¡...!).
\end{enumerate}
\end{methode}

\courssub{Extraire les arguments du texte support}

Revenons au texte étudié dans les sections précédentes sur les relations Nord-Sud. Avant de rédiger, il faut transformer nos annotations en un \textbf{tableau d'arguments}. Voici le résultat de cette étape de travail :

\begin{center}
\begin{tabularx}{\linewidth}{|X|X|X|}
\hline
\rowcolor{popblueL} \textbf{Argument du texte} & \textbf{Citation / indice} & \textbf{Exemple personnel possible} \\
\hline
La coopération doit être mutuelle, pas unilatérale & \esp{« la cooperación debe ser mutua »} & Partenariats Cameroun--Espagne dans l'éducation \\
\hline
L'aide est souvent soumise à des conditions (\esp{condicionalidad}) & \esp{« ayudas condicionadas »} & Réformes exigées par les institutions financières \\
\hline
La souveraineté des États du Sud doit être respectée & \esp{« respeto a la soberanía »} & L'Union africaine défend les décisions africaines \\
\hline
La dette freine le développement & \esp{« la deuda externa »} & Annulations de dettes en Amérique latine \\
\hline
\end{tabularx}
\end{center}

Ce tableau est votre \emph{plan de travail} : chaque ligne peut devenir une partie du développement. Au Bac, choisissez-en deux ou trois, pas plus.

\courssub{Rédiger l'introduction : modèle guidé}

L'introduction est la partie la plus codifiée. Voici un modèle à trous que vous pouvez adapter à n'importe quel texte politique :

\begin{exemplebox}[Modèle d'introduction à compléter]
\esp{En la actualidad, el tema de \textbf{[sujet]} suscita un vivo debate en la sociedad internacional. El texto propuesto, \textbf{[type de texte]} titulado « \textbf{[titre]} », aborda la cuestión de \textbf{[problématique]}. El autor defiende la tesis de que \textbf{[thèse]}. Para analizar este texto, estudiaremos en primer lugar \textbf{[axe 1]}, en segundo lugar \textbf{[axe 2]} y, por último, \textbf{[axe 3]}.}

\medskip
Traduction : « De nos jours, le thème de \textbf{[sujet]} suscite un vif débat dans la société internationale. Le texte proposé, \textbf{[type de texte]} intitulé « \textbf{[titre]} », aborde la question de \textbf{[problématique]}. L'auteur défend la thèse selon laquelle \textbf{[thèse]}. Pour analyser ce texte, nous étudierons d'abord \textbf{[axe 1]}, ensuite \textbf{[axe 2]} et enfin \textbf{[axe 3]}. »
\end{exemplebox}

Notez les formules toutes faites (\esp{en la actualidad}, \esp{suscita un vivo debate}, \esp{defiende la tesis de que}) : ce sont des \cle{connecteurs d'introduction} à mémoriser. Attention : après \esp{la tesis de que} suivi d'un fait affirmé, on utilise l'indicatif ; après un verbe de doute ou de possibilité, le subjonctif : \esp{Es posible que la ayuda sea condicionada.} « Il est possible que l'aide soit conditionnée. »

\courssub{Construire un paragraphe type}

Chaque paragraphe du développement suit un schéma en quatre temps :

\begin{enumerate}
\item \textbf{Phrase-thème} (\esp{frase temática}) : elle annonce l'idée du paragraphe, introduite par un connecteur.
\item \textbf{Argument} (\esp{argumento}) : vous développez l'idée avec votre analyse.
\item \textbf{Exemple} (\esp{ejemplo}) : une citation courte du texte ou un fait réel.
\item \textbf{Connecteur de transition} : il prépare le paragraphe suivant.
\end{enumerate}

\begin{exemplebox}[Paragraphe type]
\esp{En primer lugar, la cooperación debe ser mutua. En efecto, el autor subraya que la ayuda del Norte no puede ser un gesto de caridad, sino un intercambio entre socios iguales. Así, afirma que « la cooperación debe ser mutua », lo que demuestra su voluntad de igualdad. Sin embargo, esta igualdad no siempre se respeta en la realidad.}

\medskip
Traduction : « En premier lieu, la coopération doit être mutuelle. En effet, l'auteur souligne que l'aide du Nord ne peut être un geste de charité, mais un échange entre partenaires égaux. Ainsi, il affirme que "la coopération doit être mutuelle", ce qui démontre sa volonté d'égalité. Cependant, cette égalité n'est pas toujours respectée dans la réalité. »
\end{exemplebox}

\begin{center}
\begin{tabularx}{\linewidth}{|X|X|X|}
\hline
\rowcolor{popblueL} \textbf{Fonction} & \textbf{Español} & \textbf{Français} \\
\hline
Ajouter & \esp{además, asimismo} & de plus, également \\
\hline
Opposer & \esp{sin embargo, no obstante} & cependant, néanmoins \\
\hline
Expliquer & \esp{es decir, en efecto} & c'est-à-dire, en effet \\
\hline
Conclure & \esp{por lo tanto, en síntesis} & par conséquent, en synthèse \\
\hline
Ordonner & \esp{en primer lugar, por último} & en premier lieu, enfin \\
\hline
Donner son avis & \esp{en mi opinión, considero que} & à mon avis, je considère que \\
\hline
\end{tabularx}
\end{center}

\begin{attention}[Faux amis et calques fréquents]
\begin{itemize}
\item \esp{actualmente} signifie « actuellement, de nos jours » ; pour dire « en fait, en réalité », employez \esp{de hecho} ou \esp{en realidad}.
\item Ne traduisez pas « selon moi » par \esp{según mí} (tournure fautive) : dites \esp{en mi opinión} ou \esp{a mi juicio}.
\item « L'auteur pense que + subjonctif » est une erreur : \esp{el autor piensa que la ayuda es necesaria} (indicatif, car c'est une affirmation). Le subjonctif n'apparaît qu'avec la négation : \esp{no piensa que la ayuda sea suficiente} « il ne pense pas que l'aide soit suffisante ».
\item N'oubliez pas les points d'interrogation et d'exclamation ouvrants : ¿ ¡.
\end{itemize}
\end{attention}

\courssub{Commentaire modèle annoté}

Voici un commentaire complet d'environ 180 mots, rédigé selon la méthode. Les annotations en français identifient les procédés rhétoriques employés.

\begin{textebox}[Comentario modelo : « Las relaciones Norte-Sur : ¿cooperación o dominación ? »]
\esp{En la actualidad, las relaciones entre los países del Norte y del Sur suscitan un vivo debate. El texto propuesto aborda la cuestión de la cooperación internacional y defiende la tesis de que la ayuda debe ser mutua y respetuosa con la soberanía de los Estados.} [Introduction : accroche + sujet + thèse]

\esp{En primer lugar, el autor subraya que la cooperación no puede ser un acto de caridad, sino un intercambio entre socios iguales. Así, afirma que « la cooperación debe ser mutua », lo que revela una visión igualitaria del desarrollo.} [Partie 1 : phrase-thème + citation + analyse]

\esp{Sin embargo, denuncia también la condicionalidad de las ayudas: los países ricos imponen reformas que limitan la libertad de decisión de los países pobres. Este argumento es muy convincente, porque la deuda externa sigue frenando el desarrollo de África y América Latina.} [Partie 2 : connecteur d'opposition + évaluation personnelle]

\esp{En mi opinión, el autor tiene razón: solo una cooperación basada en el respeto permitirá un desarrollo justo.} [Opinion personnelle justifiée]

\esp{En síntesis, este texto nos invita a repensar las relaciones internacionales. Cabe preguntarse si la globalización podrá algún día beneficiar a todos los pueblos por igual.} [Conclusion : bilan + ouverture]
\end{textebox}

Glossaire : \esp{suscitar} = susciter ; \esp{socio} = partenaire ; \esp{condicionalidad} = conditionnalité ; \esp{frenar} = freiner ; \esp{cabe preguntarse} = on peut se demander.

\begin{savaistu}[Le commentaire au Bac]
Au Bac, un commentaire bien structuré peut rapporter jusqu'à 6 points sur 10 dans la grille d'expression écrite. Les correcteurs camerounais valorisent particulièrement trois choses : un plan visible (connecteurs d'ordre), des citations courtes du texte entre guillemets, et une opinion personnelle clairement signalée (\esp{en mi opinión}). Une copie sans plan apparent, même bien écrite, perd systématiquement des points de cohérence.
\end{savaistu}

\courssub{Critères d'évaluation au Bac}

Votre copie sera évaluée selon quatre critères principaux :

\begin{center}
\begin{tabularx}{\linewidth}{|X|X|}
\hline
\rowcolor{popblueL} \textbf{Critère} & \textbf{Ce que le correcteur attend} \\
\hline
Pertinence (\esp{pertinencia}) & Vous répondez au sujet ; vos analyses concernent bien le texte proposé. \\
\hline
Cohérence (\esp{coherencia}) & Plan visible, connecteurs, paragraphes articulés logiquement. \\
\hline
Richesse lexicale (\esp{riqueza léxica}) & Vocabulaire précis du champ politique : \esp{cooperación, soberanía, desarrollo, deuda, globalización}. \\
\hline
Correction grammaticale (\esp{corrección gramatical}) & Accords, conjugaisons, accents, subjonctif après les verbes d'opinion négatifs. \\
\hline
\end{tabularx}
\end{center}

\begin{exoresolu}[Transformer un résumé en commentaire]
Énoncé : un élève a écrit : \esp{« El autor dice que la cooperación debe ser mutua. Luego dice que la deuda es un problema. Después habla de la soberanía. »} Expliquez pourquoi ce passage est insuffisant et réécrivez-le en véritable paragraphe de commentaire.

\tcblower

Solution détaillée : ce passage est un \textbf{résumé} et non un commentaire : il répète ce que dit le texte (\esp{dice que... luego dice que...}) sans analyser ni évaluer. Il manque la citation exacte, l'analyse du procédé et l'opinion personnelle. Réécriture possible :

\esp{En primer lugar, el autor defiende que « la cooperación debe ser mutua », expresión que revela su rechazo de una ayuda caritativa y vertical. Además, al mencionar la deuda externa, denuncia un mecanismo que frena el desarrollo de los países del Sur. En mi opinión, este doble argumento es muy pertinente, porque vincula la cooperación con el respeto de la soberanía.}

Traduction : « En premier lieu, l'auteur défend l'idée que "la coopération doit être mutuelle", expression qui révèle son rejet d'une aide charitable et verticale. De plus, en mentionnant la dette extérieure, il dénonce un mécanisme qui freine le développement des pays du Sud. À mon avis, ce double argument est très pertinent, car il lie la coopération au respect de la souveraineté. »

On observe : un connecteur d'ordre (\esp{en primer lugar}), une citation entre guillemets, des verbes d'analyse (\esp{revela, denuncia}) et une opinion signalée (\esp{en mi opinión}).
\end{exoresolu}

\begin{exercice}[À vous de jouer]
À partir du texte support de la leçon, rédigez l'introduction complète d'un commentaire en utilisant le modèle à trous de la sous-section 4.4, puis construisez un paragraphe de développement sur l'axe « la souveraineté des États du Sud » en suivant le schéma : phrase-thème, argument, exemple, connecteur.
\end{exercice}

\begin{corrige}[Exercice]
Éléments de correction : l'introduction doit contenir les quatre éléments (accroche, sujet, thèse, plan) et au moins deux formules étudiées (\esp{en la actualidad}, \esp{defiende la tesis de que}). Le paragraphe doit s'ouvrir sur \esp{en segundo lugar} ou \esp{por otro lado}, contenir une citation courte du texte relative à \esp{la soberanía}, un verbe d'analyse (\esp{subrayar, denunciar, defender}) et se terminer par un connecteur de transition (\esp{sin embargo, por lo tanto}). Vérifiez : indicatif après \esp{defiende que}, accents (\esp{soberanía, opinión, síntesis}), et aucun calque du français.
\end{corrige}

\coursec{Production guidée : prise de position citoyenne}

Dans la section précédente, nous avons appris à \emph{commenter} un texte politique : repérer la thèse, analyser les arguments, évaluer les procédés. Il est temps maintenant de passer de l'analyse à l'\textbf{action} : vous allez à votre tour produire un texte argumentatif de type citoyen, une \esp{carta abierta} (lettre ouverte) adressée au Secrétaire général de l'ONU au sujet de la dette des pays du Sud. C'est l'aboutissement logique du module : comprendre le monde, puis y prendre la parole. Cette tâche reproduit d'ailleurs fidèlement l'épreuve d'\textbf{expression écrite} du Baccalauréat : une situation de communication, un destinataire précis, un registre imposé et un nombre de mots à respecter.

\courssub{La consigne et ses contraintes}

\begin{definition}[La tâche à accomplir]
Rédiger, en espagnol, une \cle{lettre ouverte} (\esp{carta abierta}) de \textbf{150 à 200 mots}, adressée au \esp{Secretario General de las Naciones Unidas}, pour dénoncer le poids de la \esp{deuda externa} des pays du Sud et proposer des solutions. La lettre doit obligatoirement :
\begin{itemize}
\item utiliser au moins \textbf{8 termes du lexique} travaillé dans la leçon (\esp{cooperación, desarrollo, deuda, globalización, relaciones internacionales, democracia, derechos, deberes, justicia, solidaridad}...) ;
\item respecter la \textbf{structure formelle} : en-tête (lieu et date, destinataire), formule d'appel, introduction, argumentation, conclusion, formule de politesse finale, signature ;
\item employer le \esp{usted} (jamais \esp{tú}) et des temps verbaux cohérents.
\end{itemize}
\end{definition}

\begin{attention}[Tú ou usted ?]
Dans une lettre officielle espagnole, on s'adresse au destinataire avec \esp{usted} (troisième personne du singulier), jamais avec \esp{tú}. Le tutoiement serait une faute de protocole grave. Attention aux conjugaisons : on écrit \esp{Le ruego que \textbf{considere}} et non \esp{*consideres}. De même, le possessif est \esp{su} : \esp{su atención}, \esp{su gobierno}, et non \esp{*tu atención}. Enfin, les pronoms objets se mettent à la troisième personne : \esp{le escribo}, \esp{le agradezco}, et non \esp{*te escribo}, \esp{*te agradezco}.
\end{attention}

\courssub{Observer d'abord : un modèle de lettre formelle}

Avant d'écrire, observons la structure d'une lettre formelle espagnole. Lisez ce modèle : les emplacements entre crochets sont à compléter avec vos propres éléments.

\begin{textebox}[Modelo de carta formal (a completar)]
Yaoundé, a [día] de [mes] de 2025\par
\medskip
Al Excelentísimo Señor Secretario General\par
de las Naciones Unidas\par
Nueva York\par
\medskip
\textbf{Asunto:} la deuda externa de los países del Sur\par
\medskip
Excelentísimo Señor Secretario General:\par
\medskip
Me dirijo a usted, en nombre de los jóvenes de mi país, para expresarle nuestra profunda preocupación por [tema]. En primer lugar, [primer argumento]. Además, [segundo argumento con un dato o ejemplo]. Sin embargo, [matiz o contraargumento].\par
\medskip
Por eso, le ruego que considere [primera propuesta] y que apoye [segunda propuesta]. Estamos convencidos de que la cooperación internacional y la solidaridad entre los pueblos son el único camino hacia un desarrollo justo.\par
\medskip
A la espera de su respuesta, le agradezco de antemano su atención y le ruego que acepte, Excelentísimo Señor Secretario General, la expresión de mi más alta consideración.\par
\medskip
[Nombre y apellido]\par
[Ciudad, país]
\end{textebox}

\textbf{Glossaire :} \esp{me dirijo a usted} = je m'adresse à vous ; \esp{preocupación} = inquiétude ; \esp{sin embargo} = cependant ; \esp{le ruego que considere} = je vous prie de considérer ; \esp{a la espera de su respuesta} = dans l'attente de votre réponse ; \esp{de antemano} = d'avance ; \esp{mi más alta consideración} = ma très haute considération ; \esp{estamos convencidos de que} = nous sommes convaincus que.

\begin{propriete}[La structure d'une lettre formelle espagnole]
\begin{enumerate}
\item \textbf{En-tête} : lieu et date (\esp{Yaoundé, a 15 de mayo de 2025} — notez la préposition \esp{a} devant la date, absente du français, et le mois sans majuscule).
\item \textbf{Destinataire} : titre protocolaire complet (\esp{Al Excelentísimo Señor Secretario General de las Naciones Unidas}).
\item \textbf{Objet} : \esp{Asunto:} suivi du thème.
\item \textbf{Formule d'appel} : \esp{Excelentísimo Señor...:} (avec deux-points, jamais de virgule).
\item \textbf{Corps} : introduction (qui je suis, pourquoi j'écris), arguments avec connecteurs, propositions.
\item \textbf{Formule finale} : longue et solennelle (\esp{le ruego que acepte... la expresión de mi más alta consideración}).
\item \textbf{Signature} : nom, prénom, ville.
\end{enumerate}
\end{propriete}

\begin{exemplebox}[Formules de politesse : espagnol / français]
\begin{center}
\begin{tabularx}{\linewidth}{|X|X|}
\hline
\rowcolor{popblueL} \textbf{Español} & \textbf{Français} \\
\hline
Me dirijo a usted para expresarle... & Je m'adresse à vous pour vous exprimer... \\
\hline
Le ruego que considere nuestra propuesta. & Je vous prie de considérer notre proposition. \\
\hline
Le agradezco de antemano su atención. & Je vous remercie d'avance de votre attention. \\
\hline
A la espera de su respuesta... & Dans l'attente de votre réponse... \\
\hline
Reciba la expresión de mi más alta consideración. & Veuillez agréer l'expression de ma très haute considération. \\
\hline
\end{tabularx}
\end{center}
\end{exemplebox}

\begin{attention}[Subjonctif obligatoire après « rogar que »]
Après \esp{le ruego que...} (je vous prie de...), l'espagnol exige le \textbf{subjonctif présent}, car on exprime une demande, un souhait, non un fait : \esp{Le ruego que \textbf{considere} esta situación} (je vous prie de considérer cette situation). Erreur fréquente des francophones : écrire \esp{*le ruego que considera} (indicatif) par calque du français. Rappel des formes à la troisième personne du singulier : \esp{considerar} → \esp{considere} ; \esp{apoyar} → \esp{apoye} ; \esp{cancelar} → \esp{cancele} ; \esp{hacer} → \esp{haga} ; \esp{tomar medidas} → \esp{tome medidas}. Même exigence après \esp{espero que...} : \esp{Espero que \textbf{responda} pronto} (j'espère que vous répondrez vite).
\end{attention}

\courssub{Étape 1 : le brainstorming en groupes}

La rédaction ne commence pas par la rédaction ! En groupes de quatre, vous construisez d'abord une \textbf{carte mentale} des arguments. Chaque groupe choisit un angle : économique (la dette étouffe les budgets de la santé et de l'éducation), moral (l'injustice Nord-Sud), politique (la souveraineté des États), humanitaire (les droits des populations).

\begin{popfigure}
\begin{center}
\begin{tikzpicture}[
  central/.style={draw=popblue, fill=popblueL, rounded corners, align=center, font=\bfseries, minimum width=3.2cm, minimum height=1.1cm},
  branch/.style={draw=popgreen, fill=popgreenL, rounded corners, align=center, minimum width=2.6cm, font=\small},
  lien/.style={-{Stealth[length=2.5mm]}, thick, popmuted}
]
\node[central, align=center] (c){La deuda\\de los países del Sur};
\node[branch, above left=1.2cm and 1.6cm of c, align=center] (a){Argumento\\económico};
\node[branch, below left=1.2cm and 1.6cm of c, align=center] (b){Argumento\\moral};
\node[branch, above right=1.2cm and 1.6cm of c, align=center] (d){Argumento\\político};
\node[branch, below right=1.2cm and 1.6cm of c, align=center] (e){Propuestas\\de solución};
\draw[lien] (c) -- (a);
\draw[lien] (c) -- (b);
\draw[lien] (c) -- (d);
\draw[lien] (c) -- (e);
\end{tikzpicture}
\end{center}
\legende{Carte mentale de départ : un thème central, quatre branches d'arguments à enrichir en groupe.}
\end{popfigure}

\begin{experience}[Brainstorming guidé (15 minutes)]
\begin{enumerate}
\item Par groupes de quatre, complétez la carte mentale : chaque branche doit recevoir au moins \textbf{deux arguments} et \textbf{un exemple} (un pays d'Afrique, la Guinée équatoriale, le Cameroun, l'Amérique latine...).
\item Listez ensemble les \textbf{8 mots du lexique} que vous comptez utiliser et vérifiez leur sens dans la section 3 de la leçon.
\item Choisissez les \textbf{connecteurs} de votre argumentation : \esp{en primer lugar} (d'abord), \esp{además} (de plus), \esp{sin embargo} (cependant), \esp{por eso} (c'est pourquoi), \esp{en conclusión} (en conclusion).
\item Un rapporteur par groupe présente la carte mentale à la classe en une minute.
\end{enumerate}
\end{experience}

\courssub{Étape 2 : la rédaction individuelle}

Chaque élève rédige maintenant seul sa lettre (150-200 mots), en s'appuyant sur la carte mentale du groupe et sur le modèle formel de la section 5.2. Le déroulement complet de l'activité est résumé ci-dessous.

\begin{popfigure}
\begin{center}
\begin{tikzpicture}[
  etape/.style={draw=popblue, fill=popblueL, rounded corners, align=center, minimum width=2.9cm, minimum height=1.2cm, font=\small},
  opt/.style={draw=poporange, fill=poporangeL, rounded corners, align=center, minimum width=2.9cm, minimum height=1.2cm, font=\small},
  fleche/.style={-{Stealth[length=3mm]}, thick, popdark}
]
\node[etape, align=center] (brain){Brainstorming\\(carte mentale)};
\node[etape, right=1.6cm of brain, align=center] (draft){Rédaction\\individuelle};
\node[etape, right=1.6cm of draft, align=center] (peer){Relecture\\croisée};
\node[opt, below=1.4cm of draft, align=center] (oral){Présentation\\orale (2 min)};
\draw[fleche] (brain) -- (draft);
\draw[fleche] (draft) -- (peer);
\draw[fleche] (draft) -- (oral);
\end{tikzpicture}
\end{center}
\legende{Les quatre étapes de la production guidée : de l'idée collective à la parole individuelle.}
\end{popfigure}

\begin{methode}[Checklist avant de rendre votre lettre]
\begin{enumerate}
\item \textbf{8 mots-clés} du lexique de la leçon sont-ils présents et correctement employés ? Soulignez-les.
\item \textbf{Connecteurs} : y a-t-il au moins quatre connecteurs variés (\esp{en primer lugar, además, sin embargo, por eso, en conclusión}) ?
\item \textbf{Temps verbaux cohérents} : présent de vérité générale (\esp{la deuda impide el desarrollo}), subjonctif après \esp{le ruego que...}, futur ou conditionnel pour les propositions (\esp{sería justo cancelar la deuda} — il serait juste d'annuler la dette).
\item \textbf{Formules de politesse} : appel solennel, \esp{usted} partout, formule finale complète, aucun \esp{tú}.
\item \textbf{Orthographe} : accents (\esp{cooperación, política, educación}), \esp{ñ}, points d'exclamation et d'interrogation ouvrants (¡, ¿).
\end{enumerate}
\end{methode}

\courssub{Étape 3 : la relecture croisée}

Échangez votre lettre avec un voisin. Chacun corrige la copie de l'autre à l'aide de la grille d'auto-correction ci-dessous, en notant ses remarques au crayon, avec bienveillance et précision.

\begin{center}
\begin{tabularx}{\linewidth}{|X|X|X|}
\hline
\rowcolor{popblueL} \textbf{Critère} & \textbf{À vérifier} & \textbf{Points} \\
\hline
Lexique & 8 termes du champ lexical présents et corrects & /4 \\
\hline
Structure & En-tête, appel, corps, formule finale, signature & /4 \\
\hline
Registre & \esp{usted} partout, aucun tutoiement & /3 \\
\hline
Grammaire & Subjonctif après \esp{rogar que}, accords, temps & /4 \\
\hline
Argumentation & Au moins 3 arguments + connecteurs variés & /3 \\
\hline
Orthographe & Accents, ponctuation espagnole & /2 \\
\hline
\end{tabularx}
\end{center}

\begin{exoresolu}[Corriger un extrait de lettre d'élève]
\textbf{Énoncé.} Voici l'extrait rédigé par un élève de Terminale A. Relevez et corrigez les quatre erreurs.\par
\medskip
\emph{« Hola Secretario General, te escribo porque la deuda de los países del Sur es un problema grave. Te ruego que cancelas esta deuda. Tu atención es muy importante para nosotros. »}
\tcblower
\textbf{Solution détaillée.}
\begin{enumerate}
\item \esp{Hola} : registre trop familier pour une lettre officielle. Correction : \esp{Excelentísimo Señor Secretario General:}
\item \esp{te escribo} : tutoiement interdit. Correction : \esp{le escribo} (ou mieux : \esp{me dirijo a usted}).
\item \esp{Te ruego que cancelas} : après \esp{rogar que}, il faut le subjonctif présent. Correction : \esp{Le ruego que \textbf{cancele} esta deuda.}
\item \esp{Tu atención} : possessif du tutoiement. Correction : \esp{\textbf{Su} atención es muy importante para nosotros.}
\end{enumerate}
Version corrigée : \esp{Excelentísimo Señor Secretario General: me dirijo a usted porque la deuda de los países del Sur es un problema grave. Le ruego que cancele esta deuda. Su atención es muy importante para nosotros.}
\end{exoresolu}

\courssub{Étape 4 : la présentation orale (optionnel)}

Les volontaires lisent leur lettre devant la classe en \textbf{deux minutes}, comme s'ils s'adressaient réellement à l'Assemblée générale des Nations unies. Critères d'évaluation orale : prononciation des mots du lexique (\esp{cooperación} avec l'accent tonique sur la dernière syllabe, \esp{globalización} avec le \esp{z} prononcé comme un « th » anglais en Espagne ou comme un « s » en Amérique latine), intonation solennelle, regard vers l'auditoire, respect du temps.

\begin{savaistu}[Le protocole diplomatique des lettres internationales]
Les lettres adressées aux organisations internationales suivent un protocole strict : on écrit dans l'une des \textbf{langues officielles} de l'organisation (l'ONU en compte six : espagnol, anglais, français, arabe, chinois et russe), on utilise le titre exact du destinataire (\esp{Excelentísimo Señor Secretario General}) et une formule finale de « haute considération ». L'espagnol est ainsi l'une des grandes langues de la diplomatie mondiale : le maîtriser, c'est pouvoir porter la voix du Cameroun et de l'Afrique dans les instances internationales. La Guinée équatoriale, voisine du Cameroun, est d'ailleurs le seul pays d'Afrique subsaharienne dont l'espagnol est langue officielle.
\end{savaistu}

\begin{exercice}[À vous d'écrire]
Rédigez votre lettre ouverte au Secrétaire général de l'ONU (150-200 mots) en suivant toutes les étapes de cette section : carte mentale en groupe, rédaction individuelle avec le modèle formel, vérification avec la checklist, relecture croisée avec la grille. N'oubliez pas : 8 mots du lexique, \esp{usted}, subjonctif après \esp{le ruego que}, et une conclusion qui propose au moins une solution concrète (annulation de la dette, coopération renforcée, investissement dans l'éducation).
\end{exercice}

\begin{corrige}[Exercice « À vous d'écrire » : éléments de rédaction]
Il n'existe pas une seule bonne lettre, mais voici un \textbf{modèle de réponse possible} (environ 170 mots) que l'enseignant pourra projeter après la relecture croisée :\par
\medskip
\esp{Yaoundé, a 15 de mayo de 2025. Excelentísimo Señor Secretario General: me dirijo a usted, en nombre de los jóvenes de Camerún, para expresarle nuestra preocupación por la deuda externa de los países del Sur. En primer lugar, esta deuda impide el desarrollo de nuestros pueblos: los gobiernos gastan más dinero en pagarla que en la educación y la salud. Además, la globalización ha aumentado las desigualdades entre el Norte y el Sur. Sin embargo, creemos que la cooperación internacional puede cambiar esta situación. Por eso, le ruego que considere la cancelación de la deuda y que apoye una cooperación más justa y solidaria. Los derechos de los pueblos del Sur son también deberes de la comunidad internacional. A la espera de su respuesta, le agradezco de antemano su atención y le ruego que acepte, Excelentísimo Señor Secretario General, la expresión de mi más alta consideración.}\par
\medskip
« Yaoundé, le 15 mai 2025. Très excellent Monsieur le Secrétaire général : je m'adresse à vous, au nom des jeunes du Cameroun, pour vous exprimer notre inquiétude face à la dette extérieure des pays du Sud. D'abord, cette dette empêche le développement de nos peuples : les gouvernements dépensent plus d'argent à la rembourser qu'à l'éducation et à la santé. De plus, la mondialisation a accru les inégalités entre le Nord et le Sud. Cependant, nous croyons que la coopération internationale peut changer cette situation. C'est pourquoi je vous prie de considérer l'annulation de la dette et de soutenir une coopération plus juste et solidaire. Les droits des peuples du Sud sont aussi des devoirs de la communauté internationale. Dans l'attente de votre réponse, je vous remercie d'avance de votre attention et vous prie d'agréer, Très excellent Monsieur le Secrétaire général, l'expression de ma très haute considération. »\par
\medskip
\textbf{Points de méthode à vérifier sur ce modèle :} 10 mots du lexique (\esp{deuda externa, países del Sur, desarrollo, educación, globalización, cooperación internacional, derechos, deberes, solidaria, justa}) ; deux subjonctifs corrects (\esp{considere}, \esp{apoye}, \esp{acepte}) ; \esp{usted} implicite partout ; quatre connecteurs (\esp{en primer lugar, además, sin embargo, por eso}).
\end{corrige}

\begin{aretenir}[L'essentiel de la production guidée]
Écrire une lettre officielle en espagnol, c'est conjuguer trois compétences : la \textbf{forme} (structure protocolaire, formules de politesse, \esp{usted}), la \textbf{langue} (lexique des relations internationales, subjonctif après les verbes de demande, connecteurs argumentatifs) et la \textbf{pensée citoyenne} (des arguments clairs, des exemples précis, des propositions concrètes). La démarche en quatre étapes — brainstorming, rédaction, relecture croisée, présentation orale — transforme un exercice scolaire en véritable acte de communication internationale, exactement ce que l'épreuve d'expression écrite du Baccalauréat attend de vous.
\end{aretenir}

\newpage

\coursec{Activité d'intégration}

\courssub{Objectifs de l'activité}

À l'issue de ce débat citoyen, l'élève sera capable de :
\begin{itemize}
\item mobiliser le \cle{lexique des relations internationales} et des systèmes politiques dans un échange oral authentique ;
\item construire une \cle{argumentation structurée} (thèse, arguments, exemples, concession, conclusion) à l'oral comme à l'écrit ;
\item utiliser correctement les \cle{connecteurs logiques} et le \cle{subjonctif} après les verbes d'opinion et de doute ;
\item écouter, réagir et réfuter poliment le point de vue d'un adversaire ;
\item rédiger collectivement une synthèse courte et équilibrée.
\end{itemize}

\courssub{Situation}

Dans le cadre de la Semaine de la Francophonie et des Langues, le club d'espagnol du lycée de Yaoundé organise un grand débat public devant les classes de Première et de Terminale. Le thème retenu cette année porte sur la coopération entre les pays riches et les pays en développement. Le proviseur a demandé aux élèves de Terminale A de préparer ce débat en s'appuyant sur la presse. La veille, la journaliste du club a distribué le texte suivant, extrait d'un article original rédigé pour l'occasion :

\begin{textebox}[La cooperación Norte-Sur, ¿un modelo agotado?]
Desde hace más de sesenta años, la cooperación entre los países del Norte y los del Sur ha movilizado miles de millones de dólares en ayuda al desarrollo. Sin embargo, los resultados siguen siendo discutidos. Muchos países africanos, como Camerún, han diversificado sus socios: ya no dependen únicamente de sus antiguas potencias coloniales, sino que cooperan también con China, con los países del Golfo o con sus vecinos africanos dentro de la Unión Africana.

Para algunos analistas, la ayuda tradicional crea dependencia y no resuelve los problemas de fondo: la deuda externa, los precios bajos de las materias primas como el cacao o el café, y la fuga de cerebros. Otros, en cambio, recuerdan que la cooperación ha permitido construir carreteras, hospitales y escuelas, y que organismos como la ONU desempeñan un papel esencial en la paz y los derechos humanos.

La verdadera cuestión quizá no sea si la cooperación debe continuar, sino cómo transformarla: menos caridad y más comercio justo, menos condiciones impuestas y más respeto mutuo. En un mundo globalizado, los destinos del Norte y del Sur están más unidos que nunca.
\end{textebox}

\textbf{Glossaire :} \esp{agotado} = épuisé ; \esp{ayuda al desarrollo} = aide au développement ; \esp{discutido} = contesté, discuté ; \esp{socios} = partenaires ; \esp{potencias coloniales} = puissances coloniales ; \esp{fuga de cerebros} = fuite des cerveaux ; \esp{deuda externa} = dette extérieure ; \esp{materias primas} = matières premières ; \esp{carreteras} = routes ; \esp{desempeñar un papel} = jouer un rôle ; \esp{comercio justo} = commerce équitable ; \esp{respeto mutuo} = respect mutuel.

\textbf{Traduction du texte :} « Depuis plus de soixante ans, la coopération entre les pays du Nord et ceux du Sud a mobilisé des milliards de dollars d'aide au développement. Cependant, les résultats restent contestés. De nombreux pays africains, comme le Cameroun, ont diversifié leurs partenaires : ils ne dépendent plus uniquement de leurs anciennes puissances coloniales, mais coopèrent aussi avec la Chine, avec les pays du Golfe ou avec leurs voisins africains au sein de l'Union africaine. Pour certains analystes, l'aide traditionnelle crée une dépendance et ne résout pas les problèmes de fond : la dette extérieure, les prix bas des matières premières comme le cacao ou le café, et la fuite des cerveaux. D'autres, en revanche, rappellent que la coopération a permis de construire des routes, des hôpitaux et des écoles, et que des organismes comme l'ONU jouent un rôle essentiel dans la paix et les droits de l'homme. La véritable question n'est peut-être pas de savoir si la coopération doit continuer, mais comment la transformer : moins de charité et plus de commerce équitable, moins de conditions imposées et plus de respect mutuel. Dans un monde globalisé, les destins du Nord et du Sud sont plus liés que jamais. »

Tu es membre du club d'espagnol. Tu dois prendre position, avec ton camp, sur la question : \emph{¿La cooperación Norte-Sur es todavía pertinente?}

\courssub{Consignes}

\begin{enumerate}
\item \textbf{Lecture active.} Relis le texte et souligne en rouge trois arguments \emph{contre} la coopération traditionnelle, en vert trois arguments \emph{pour}. Note dans la marge les mots du lexique de la leçon qui reviennent (\esp{cooperación}, \esp{deuda}, \esp{desarrollo}, \esp{globalización}, \esp{ONU}, \esp{Unión Africana}).

\item \textbf{Préparation en groupe (10 minutes).} La classe est divisée en deux camps : le groupe A défend la pertinence de la coopération Nord-Sud, le groupe B la conteste. Chaque groupe rédige une fiche avec : une \cle{thèse} claire (une phrase), \cle{trois arguments} hiérarchisés, un \cle{exemple} concret pour chaque argument (le texte, le Cameroun, l'Afrique, le monde hispanophone), et une \cle{concession} possible (\esp{Aunque es cierto que...}).

\item \textbf{Choix des porte-parole.} Chaque groupe désigne deux orateurs et un secrétaire. Un élève volontaire joue le rôle de \cle{modérateur} : il donne la parole, veille au respect du temps et rappelle les règles de courtoisie.

\item \textbf{Débat (15 minutes).} Alternance des prises de parole : deux minutes maximum par intervention. Chaque orateur doit employer au moins \cle{trois connecteurs logiques} différents et \cle{une phrase au subjonctif} (par exemple après \esp{no creo que}, \esp{es posible que}, \esp{aunque} suivi d'une hypothèse).

\item \textbf{Réfutation.} Après le premier tour, chaque camp dispose de trois minutes pour répondre aux arguments adverses, en commençant obligatoirement par une formule de concession polie.

\item \textbf{Vote.} Le public (les autres élèves) vote à main levée pour le camp le plus convaincant. Le modérateur proclame le résultat.

\item \textbf{Synthèse écrite collective.} Les deux secrétaires rédigent ensemble, au tableau, une synthèse de cinq lignes environ qui présente les deux positions et propose une conclusion équilibrée. La classe corrige et améliore le texte avant de l'afficher.
\end{enumerate}

\begin{attention}[Règles d'or du débat]
On attaque les \emph{idées}, jamais les personnes. On attend son tour de parole. Toute affirmation doit être appuyée par un argument ou un exemple. Interdit de lire intégralement ses notes : on parle en regardant l'auditoire.
\end{attention}

\courssub{Ressources à mobiliser}

\begin{aretenir}[Lexique utile pour le débat]
\begin{itemize}
\item \textbf{La coopération :} \esp{la ayuda al desarrollo}, \esp{los socios}, \esp{el intercambio}, \esp{el comercio justo}, \esp{la dependencia}, \esp{la solidaridad internacional}.
\item \textbf{Les institutions :} \esp{la ONU}, \esp{la Unión Africana}, \esp{los organismos internacionales}, \esp{la comunidad internacional}.
\item \textbf{Les systèmes politiques :} \esp{la democracia}, \esp{la dictadura}, \esp{la monarquía}, \esp{la república}, \esp{las elecciones libres}.
\item \textbf{Les problèmes :} \esp{la deuda externa}, \esp{la pobreza}, \esp{la desigualdad}, \esp{la fuga de cerebros}, \esp{las materias primas}.
\item \textbf{Les droits et devoirs :} \esp{los derechos humanos}, \esp{la libertad de expresión}, \esp{el derecho a la educación}, \esp{los deberes del ciudadano}.
\end{itemize}
\end{aretenir}

\begin{propriete}[Connecteurs et structures argumentatives]
\begin{itemize}
\item \textbf{Introduire son opinion :} \esp{En mi opinión...}, \esp{Desde mi punto de vista...}, \esp{Estoy convencido de que...} (suivi de l'\cle{indicatif}, car l'opinion est affirmée).
\item \textbf{Exprimer le doute ou la négation d'une opinion :} \esp{No creo que} + \cle{subjonctif} : \esp{No creo que la ayuda tradicional sea suficiente.} « Je ne crois pas que l'aide traditionnelle soit suffisante. »
\item \textbf{Ordonner :} \esp{En primer lugar}, \esp{además}, \esp{por otro lado}, \esp{finalmente}.
\item \textbf{Concéder :} \esp{Aunque es verdad que...} + indicatif (fait admis), \esp{Es cierto que..., pero...}, \esp{Sin embargo}.
\item \textbf{Conclure :} \esp{En conclusión}, \esp{Por todo ello}, \esp{En resumen}.
\end{itemize}
\end{propriete}

\begin{attention}[Pièges fréquents des francophones]
\begin{itemize}
\item \esp{Actualmente} signifie « actuellement, à l'heure actuelle » : c'est un bon équivalent. Mais attention au sens de l'adjectif \esp{actual} : il signifie « actuel, présent », jamais « réel, véritable ». Pour « en fait, en réalité », on dit \esp{en realidad} ou \esp{de hecho}.
\item \esp{No creo que} exige le subjonctif : \esp{No creo que sea justo}, et non \esp{no creo que es justo}. En revanche, à la forme affirmative, \esp{Creo que} est suivi de l'indicatif : \esp{Creo que es justo.}
\item \esp{la ONU} est féminin (la Organización de las Naciones Unidas) ; on écrit \esp{la ONU}, sans élision de l'article.
\item « Être d'accord » se dit \esp{estar de acuerdo} (avec \esp{estar}), jamais \esp{ser de acuerdo} : \esp{Estoy de acuerdo contigo.} « Je suis d'accord avec toi. »
\item « Dépendre de » se traduit par \esp{depender de} : \esp{Camerún depende de sus socios.} Ne dis pas \esp{depender en}.
\item \esp{los derechos humanos} : ne traduis pas « droits humains » mot à mot par \esp{derechos humanos} au féminin ; \esp{derecho} est masculin.
\end{itemize}
\end{attention}

\courssub{Proposition de production et corrigé commenté}

\begin{exoresolu}[Intervention modèle du camp « pour » et synthèse collective]
\textbf{Énoncé.} Rédige l'intervention d'ouverture du camp favorable à la coopération Nord-Sud (environ dix lignes), puis la synthèse collective de cinq lignes.
\tcblower
\textbf{Intervention modèle :}

\esp{Señor moderador, queridos compañeros: en primer lugar, estoy convencido de que la cooperación Norte-Sur sigue siendo necesaria. Es cierto que la ayuda tradicional tiene defectos, pero sin ella, muchos hospitales y escuelas de nuestro país no existirían. Además, organismos como la ONU defienden la paz y los derechos humanos en todo el mundo. Por otro lado, no creo que el comercio justo pueda reemplazar totalmente la cooperación: los dos deben avanzar juntos. En un mundo globalizado, los problemas del Sur son también los problemas del Norte. En conclusión, no debemos abandonar la cooperación, sino transformarla con más respeto mutuo.}

\textbf{Traduction :} « Monsieur le modérateur, chers camarades : en premier lieu, je suis convaincu que la coopération Nord-Sud reste nécessaire. Il est vrai que l'aide traditionnelle a des défauts, mais sans elle, beaucoup d'hôpitaux et d'écoles de notre pays n'existeraient pas. De plus, des organismes comme l'ONU défendent la paix et les droits de l'homme dans le monde entier. D'autre part, je ne crois pas que le commerce équitable puisse remplacer totalement la coopération : les deux doivent avancer ensemble. Dans un monde globalisé, les problèmes du Sud sont aussi les problèmes du Nord. En conclusion, nous ne devons pas abandonner la coopération, mais la transformer avec plus de respect mutuel. »

\textbf{Synthèse collective modèle (5 lignes) :}

\esp{Los dos grupos han defendido sus ideas con argumentos sólidos. Para unos, la cooperación crea dependencia y mantiene la deuda externa; para otros, ha construido carreteras, escuelas y hospitales. Todos estamos de acuerdo en que el modelo debe cambiar: menos caridad y más comercio justo. La cooperación sigue siendo pertinente si respeta la dignidad de los pueblos del Sur.}

\textbf{Traduction :} « Les deux groupes ont défendu leurs idées avec des arguments solides. Pour les uns, la coopération crée une dépendance et entretient la dette extérieure ; pour les autres, elle a construit des routes, des écoles et des hôpitaux. Tous, nous sommes d'accord pour dire que le modèle doit changer : moins de charité et plus de commerce équitable. La coopération reste pertinente si elle respecte la dignité des peuples du Sud. »

\textbf{Commentaire des choix de langue :}
\begin{itemize}
\item \esp{estoy convencido de que} + indicatif (\esp{sigue siendo}) : opinion affirmée, donc pas de subjonctif ;
\item \esp{no creo que... pueda} : subjonctif correctement employé après un verbe d'opinion à la forme négative ;
\item connecteurs variés : \esp{en primer lugar}, \esp{además}, \esp{por otro lado}, \esp{en conclusión} ;
\item concession honnête : \esp{Es cierto que... pero...}, ce qui renforce la crédibilité de l'orateur ;
\item chute efficace : \esp{no debemos abandonar..., sino transformarla}, structure de correction \esp{no... sino...} très valorisée à l'oral ;
\item conditionnel bien formé dans \esp{no existirían} pour exprimer l'irréel du présent (« n'existeraient pas »).
\end{itemize}
\end{exoresolu}

\courssub{Bilan de l'activité}

Cette activité d'intégration a permis de réunir toutes les compétences de la leçon dans une tâche finale authentique. La \cle{compréhension de l'écrit} a été mobilisée pour extraire du texte les arguments des deux camps ; le \cle{lexique des relations internationales} et des systèmes politiques a été réemployé dans un contexte oral réel ; la \cle{grammaire} (subjonctif, connecteurs, structures d'opinion) a servi l'argumentation et non l'inverse. Le débat a en outre développé des compétences transversales essentielles au profil de sortie du lycéen camerounais : écoute active, respect de la parole d'autrui, esprit critique et prise de position citoyenne. La synthèse écrite collective, affichée au tableau, constitue une trace écrite réutilisable pour l'épreuve de commentaire et d'expression du baccalauréat. Pour aller plus loin, la classe pourra rédiger un article pour le journal du lycée relatant le débat et son résultat, ou préparer un second débat sur les droits et devoirs du citoyen dans une démocratie.

\newpage

\coursec{Méthodes et exercices résolus}

\courssub{Fiche méthode 1 : Lire et comprendre un texte sur les relations internationales}

\begin{textebox}[Texto de apoyo : « Un comercio más justo »]
\esp{Las relaciones internacionales entre los países del Norte y los del Sur siguen marcadas por una gran desigualdad. Los países ricos compran las materias primas del Sur — el cacao, el café, el petróleo — a precios muy bajos y venden después productos manufacturados mucho más caros. Además, la deuda externa impide a muchos Estados africanos invertir en la educación y en la sanidad. Sin embargo, la Unión Africana y la ONU defienden una cooperación basada en el comercio justo y no solamente en la ayuda humanitaria. «No queremos limosna, queremos justicia», declaró un dirigente africano durante la última cumbre. Por lo tanto, cada vez más ciudadanos, en el Norte como en el Sur, exigen unas relaciones internacionales más equitativas.}
\end{textebox}

\textbf{Glossaire du texte :} \esp{las materias primas} = les matières premières ; \esp{la deuda externa} = la dette extérieure ; \esp{la sanidad} = la santé (système sanitaire) ; \esp{la limosna} = l'aumône ; \esp{un dirigente} = un dirigeant ; \esp{la cumbre} = le sommet (réunion internationale) ; \esp{equitativo} = équitable.

\begin{methode}[Comprendre un texte de presse sur les relations Nord-Sud]
\begin{enumerate}
\item \textbf{Première lecture globale.} Lis le texte entier sans dictionnaire. Repère le \cle{tema} (sujet), le \cle{tipo de texto} (article, éditorial, discours) et la \cle{fuente} (source). Souligne les mots transparents : \esp{cooperación}, \esp{globalización}, \esp{desarrollo}, \esp{deuda}.
\item \textbf{Repérage des acteurs.} Entoure les noms propres et sigles : \esp{la ONU} (l'ONU), \esp{la Unión Africana}, \esp{los países del Norte / del Sur}. Note qui fait quoi : \esp{donar} (donner), \esp{prestar} (prêter), \esp{invertir} (investir), \esp{exigir} (exiger), \esp{reclamar} (réclamer).
\item \textbf{Recherche de la thèse.} Un texte sur les relations Nord-Sud oppose souvent deux points de vue : aide ou échange équitable ? Annulation de la dette ou remboursement ? Repère les connecteurs : \esp{sin embargo} (cependant), \esp{por lo tanto} (par conséquent), \esp{no obstante} (néanmoins), \esp{además} (de plus).
\item \textbf{Lexique en contexte.} Devine le sens des mots inconnus par le contexte avant de chercher : \esp{el endeudamiento} vient de \esp{deuda} (dette) $\rightarrow$ « endettement ».
\item \textbf{Réponse aux questions.} Réponds toujours par une \textbf{phrase complète en espagnol}, en reprenant les termes de la question, et cite le texte entre guillemets pour justifier : \esp{Según el texto, «...»}.
\end{enumerate}
\end{methode}

\begin{exoresolu}[Questions de compréhension sur un article]
\textbf{Texte :} \esp{La cooperación entre los países del Norte y los del Sur sigue siendo desigual. Según los analistas, la deuda externa de muchos países africanos impide el desarrollo de la educación y de la sanidad. Sin embargo, la Unión Africana reclama un comercio más justo y no solamente ayuda humanitaria.}

\textbf{Questions :} 1) \esp{¿Qué problema principal denuncia el texto?} 2) \esp{¿Qué consecuencias tiene la deuda externa?} 3) \esp{¿Qué reclama la Unión Africana?} 4) Trouve dans le texte un connecteur d'opposition et traduis-le.
\tcblower
\textbf{1)} \esp{El texto denuncia que la cooperación entre los países del Norte y los del Sur sigue siendo desigual.} — On reprend la formulation de la question et on cite le texte : « \esp{sigue siendo desigual} ».

\textbf{2)} \esp{La deuda externa impide el desarrollo de la educación y de la sanidad.} — Justification : « \esp{impide el desarrollo...} ». Attention : \esp{impide} vient du verbe \esp{impedir} (empêcher), verbe à changement de voyelle e $\rightarrow$ i.

\textbf{3)} \esp{La Unión Africana reclama un comercio más justo y no solamente ayuda humanitaria.} — \esp{reclamar} = réclamer. Il faut lire la phrase jusqu'au bout : l'Union africaine ne réclame pas de l'aide, elle réclame un commerce plus juste.

\textbf{4)} Le connecteur d'opposition est \esp{sin embargo}, qui signifie « cependant ». Il introduit la nuance : les pays du Sud ne veulent pas seulement recevoir, ils veulent échanger équitablement.

\textbf{Conclusion :} chaque réponse est une phrase complète en espagnol, appuyée sur une citation exacte du texte : c'est ce que le correcteur attend.
\end{exoresolu}

\courssub{Fiche méthode 2 : Maîtriser le lexique des relations internationales et des systèmes politiques}

\begin{methode}[Construire et réutiliser le lexique thématique]
\begin{enumerate}
\item \textbf{Classe par champs lexicaux.} Trois champs pour cette leçon : (a) relations internationales : \esp{la cooperación}, \esp{la ayuda}, \esp{el comercio}, \esp{la globalización}, \esp{la deuda externa}, \esp{el desarrollo} ; (b) institutions : \esp{la ONU}, \esp{la Unión Africana}, \esp{el gobierno}, \esp{el parlamento} ; (c) systèmes politiques : \esp{la democracia}, \esp{la dictadura}, \esp{la monarquía}, \esp{la república}.
\item \textbf{Apprends les familles de mots.} \esp{desarrollar} (verbe) $\rightarrow$ \esp{el desarrollo} (nom) $\rightarrow$ \esp{desarrollado / en vías de desarrollo} (adjectifs). \esp{gobernar} $\rightarrow$ \esp{el gobierno} $\rightarrow$ \esp{gobernante}.
\item \textbf{Mémorise les collocations.} On dit \esp{anular la deuda} (annuler la dette), \esp{respetar los derechos humanos}, \esp{celebrar elecciones} (organiser des élections), \esp{ejercer el derecho al voto}.
\item \textbf{Réutilise dans une phrase.} Un mot n'est acquis que s'il est employé : fabrique une phrase par mot nouveau, par exemple \esp{Los ciudadanos tienen el derecho de votar y el deber de respetar las leyes.}
\item \textbf{Vérifie le genre.} \esp{la deuda}, \esp{la ONU} (féminin : \esp{la Organización de las Naciones Unidas}), \esp{el desarrollo}, \esp{el derecho}.
\end{enumerate}
\end{methode}

\begin{center}
\begin{tabularx}{\linewidth}{|X|X|X|}
\hline
\rowcolor{popblueL} Español & Français & Remarque \\
\hline
\esp{la democracia} & la démocratie & le peuple élit ses représentants \\
\hline
\esp{la dictadura} & la dictature & un seul « c » en espagnol \\
\hline
\esp{la monarquía} & la monarchie & accent sur le í ; \esp{el rey} = le roi \\
\hline
\esp{la república} & la république & \esp{el presidente} = le président \\
\hline
\esp{la deuda externa} & la dette extérieure & féminin ; ne pas confondre avec \esp{la duda} (le doute) \\
\hline
\esp{los países en vías de desarrollo} & les pays en voie de développement & expression figée \\
\hline
\esp{los derechos humanos} & les droits de l'homme & jamais \esp{derechos de los hombres} \\
\hline
\esp{la libertad de expresión} & la liberté d'expression & préposition \esp{de} \\
\hline
\end{tabularx}
\end{center}

\begin{exoresolu}[Lexique : compléter et classer]
\textbf{Énoncé :} 1) Complète avec le mot qui convient : \esp{democracia, deuda, cooperación, dictadura, derechos}. a) \esp{En una \_\_\_ los ciudadanos eligen a sus representantes.} b) \esp{La \_\_\_ externa ahoga la economía de algunos países.} c) \esp{La \_\_\_ internacional es indispensable para el desarrollo.} d) \esp{En una \_\_\_ no hay libertad de expresión.} e) \esp{Los \_\_\_ humanos son universales.} 2) Traduis en français : \esp{los países en vías de desarrollo}.
\tcblower
\textbf{1) a)} \esp{democracia} — en démocratie, le pouvoir appartient au peuple qui élit ses représentants. \textbf{b)} \esp{deuda} — \esp{la deuda externa} = la dette extérieure ; attention au genre féminin. \textbf{c)} \esp{cooperación} — la coopération internationale ; mot transparent mais avec un seul « p » en espagnol. \textbf{d)} \esp{dictadura} — la dictature supprime les libertés ; contraste logique avec la phrase a). \textbf{e)} \esp{derechos} — \esp{los derechos humanos} = les droits de l'homme ; accord pluriel avec \esp{humanos}.

\textbf{2)} \esp{Los países en vías de desarrollo} = « les pays en voie de développement ». Ne jamais traduire par « pays en train de se développer » : l'expression figée française est « en voie de ».

\textbf{Conclusion :} le lexique politique est très transparent pour un francophone ; le vrai danger est l'orthographe (\esp{cooperación} sans double p, \esp{dictadura} avec un seul c) et le genre (\esp{la deuda}).
\end{exoresolu}

\courssub{Fiche méthode 3 : Commenter un texte politique}

\begin{methode}[La méthode du comentario de texto en quatre étapes]
\begin{enumerate}
\item \textbf{Introducción.} Présente le texte : type, thème, contexte. Formule : \esp{Este texto es un artículo de prensa que trata de las relaciones Norte-Sur.}
\item \textbf{Resumen.} Résume les idées principales en deux ou trois phrases, sans citer mot à mot : \esp{El autor explica que... y denuncia que...}
\item \textbf{Análisis.} Analyse la structure et les arguments : repère la thèse (\esp{la tesis}), les arguments (\esp{los argumentos}) et les procédés (connecteurs, exemples, chiffres). Formules utiles : \esp{En primer lugar... Además... Sin embargo... En conclusión...}
\item \textbf{Opinión personal.} Donne ton avis argumenté : \esp{En mi opinión... / Estoy de acuerdo con el autor porque... / Desde mi punto de vista...} Termine par une ouverture : \esp{Este problema nos concierne a todos porque...}
\end{enumerate}
\textbf{Réflexes :} rédige entièrement en espagnol ; utilise le présent de l'indicatif pour l'analyse ; évite de recopier le texte ; relie les parties avec des connecteurs.
\end{methode}

\begin{exoresolu}[Commentaire modèle sur un texte politique]
\textbf{Énoncé :} Rédige un commentaire (10-12 lignes) du texte suivant : \esp{La globalización ha acercado a los pueblos, pero también ha aumentado las diferencias entre ricos y pobres. Los países del Sur exportan materias primas baratas e importan productos caros. Por lo tanto, muchos ciudadanos exigen unas relaciones internacionales más justas.}

\tcblower
\textbf{Raisonnement :} le texte est un court éditorial ; thèse : la mondialisation crée des inégalités ; structure : concession (\esp{pero}), exemple (matières premières), conclusion (\esp{por lo tanto}).

\textbf{Comentario redactado :}

\esp{Este texto es un artículo de opinión que trata de la globalización y de las relaciones Norte-Sur. El autor explica que la globalización ha acercado a los pueblos, pero que también ha aumentado las diferencias entre ricos y pobres. Para defender su tesis, da un ejemplo concreto: los países del Sur exportan materias primas baratas e importan productos caros. El conector «por lo tanto» introduce la conclusión: los ciudadanos exigen unas relaciones internacionales más justas. En mi opinión, el autor tiene razón porque el comercio actual no beneficia a todos por igual. Por ejemplo, los países productores de cacao deberían recibir un precio más justo. En conclusión, este texto nos invita a reflexionar sobre un mundo más solidario.}

\textbf{Justifications grammaticales :} \esp{ha acercado / ha aumentado} : parfait (pretérito perfecto) pour un bilan récent avec effet présent ; \esp{deberían} : conditionnel pour exprimer un souhait (« devraient ») ; \esp{tiene razón} : expression figée, sans article ; \esp{e importan} : \esp{y} devient \esp{e} devant un mot commençant par i.

\textbf{Conclusion :} le commentaire suit les quatre étapes (présentation, résumé, analyse, opinion) et reste entièrement en espagnol.
\end{exoresolu}

\courssub{Fiche méthode 4 : Présenter droits, devoirs et libertés du citoyen}

\begin{methode}[Parler des droits et devoirs avec les bonnes structures]
\begin{enumerate}
\item \textbf{Les structures-clés.} \esp{Tener derecho a + infinitif / nom} : \esp{Los ciudadanos tienen derecho a votar.} \esp{Tener el deber de + infinitif} : \esp{Tenemos el deber de respetar las leyes.} \esp{Es obligatorio / está prohibido + infinitif} : \esp{Está prohibido discriminar.}
\item \textbf{Le vocabulaire des libertés.} \esp{la libertad de expresión}, \esp{la libertad de prensa}, \esp{la libertad de reunión}, \esp{el derecho al voto}, \esp{el derecho a la educación}, \esp{la igualdad ante la ley}.
\item \textbf{La préposition exacte.} \esp{derecho a} (jamais \esp{derecho de} pour un droit à quelque chose), \esp{deber de}, \esp{libertad de}. C'est la préposition qui porte le sens.
\item \textbf{Exprimer l'opinion citoyenne.} \esp{Pienso que + indicatif} ; \esp{No creo que + subjonctif} : \esp{No creo que la dictadura respete los derechos humanos.}
\item \textbf{Conclure par un engagement.} \esp{Debemos participar en la vida política de nuestro país.}
\end{enumerate}
\end{methode}

\begin{exoresolu}[Thème : droits et devoirs du citoyen]
\textbf{Énoncé :} Traduis en espagnol : 1) « Les citoyens ont le droit de voter et le devoir de respecter les lois. » 2) « Dans une démocratie, la liberté d'expression est un droit fondamental. » 3) « Je ne crois pas que la dictature garantisse les droits de l'homme. »
\tcblower
\textbf{1)} \esp{Los ciudadanos tienen derecho a votar y el deber de respetar las leyes.} — Structure \esp{tener derecho a + infinitivo} et \esp{el deber de + infinitivo} ; \esp{las leyes} au pluriel comme en français.

\textbf{2)} \esp{En una democracia, la libertad de expresión es un derecho fundamental.} — \esp{libertad de + nom} ; l'adjectif \esp{fundamental} est invariable en genre et se place après le nom.

\textbf{3)} \esp{No creo que la dictadura garantice los derechos humanos.} — Après \esp{no creer que}, le subjonctif est obligatoire : \esp{garantice} (subjonctif présent de \esp{garantizar}, avec changement orthographique z $\rightarrow$ c devant e). Dire « \esp{los derechos de los hombres} » serait un calque : l'expression espagnole est \esp{los derechos humanos}.

\textbf{Conclusion :} trois pièges évités : la préposition (\esp{derecho a}), le subjonctif après une opinion niée, et l'expression figée \esp{derechos humanos}.
\end{exoresolu}

\courssub{Fiche méthode 5 : Version et thème sur les relations internationales}

\begin{methode}[Traduire un passage sur la coopération Nord-Sud]
\begin{enumerate}
\item \textbf{Version (espagnol $\rightarrow$ français).} Identifie d'abord le temps dominant (souvent le présent ou le parfait). Traduis les faux amis avec prudence : \esp{actualmente} = « actuellement » (et non « effectivement »), \esp{los países ricos} = « les pays riches ».
\item \textbf{Thème (français $\rightarrow$ espagnol).} Repère les structures à calquer sur l'espagnol, pas sur le français : « il faut » $\rightarrow$ \esp{hay que + infinitif} ; « de plus en plus » $\rightarrow$ \esp{cada vez más} ; « aider à » $\rightarrow$ \esp{ayudar a + infinitif}.
\item \textbf{Accords.} Vérifie chaque adjectif : \esp{las relaciones internacionales}, \esp{una cooperación equitativa}.
\item \textbf{Accents.} \esp{cooperación}, \esp{globalización}, \esp{política}, \esp{económico}, \esp{países} (accent sur le í : \esp{país} $\rightarrow$ \esp{países}).
\item \textbf{Relecture.} Relis en ne regardant que l'espagnol : la phrase est-elle naturelle pour un hispanophone ?
\end{enumerate}
\end{methode}

\begin{exoresolu}[Thème : la coopération internationale]
\textbf{Énoncé :} Traduis en espagnol : « Il faut renforcer la coopération entre les pays du Nord et les pays du Sud. Les nations riches doivent annuler la dette des pays pauvres et investir dans l'éducation. De plus en plus de jeunes Africains exigent un commerce équitable. »
\tcblower
\textbf{Traduction :} \esp{Hay que reforzar la cooperación entre los países del Norte y los países del Sur. Las naciones ricas deben anular la deuda de los países pobres e invertir en la educación. Cada vez más jóvenes africanos exigen un comercio justo.}

\textbf{Justifications :} « il faut » impersonnel $\rightarrow$ \esp{hay que + infinitivo} (et non \esp{es necesario que}, qui demanderait le subjonctif) ; \esp{deben} : présent de \esp{deber} pour « doivent » ; \esp{invertir en} : préposition \esp{en} obligatoire ; \esp{e invertir} : \esp{y} devient \esp{e} devant i ; « de plus en plus de » $\rightarrow$ \esp{cada vez más} + nom sans article ; \esp{jóvenes africanos} : accord pluriel, adjectif de nationalité sans majuscule en espagnol ; \esp{comercio justo} : « équitable » se dit \esp{justo} dans cette collocation.

\textbf{Conclusion :} la traduction respecte les collocations espagnoles (\esp{anular la deuda}, \esp{invertir en}, \esp{comercio justo}) plutôt que de calquer le français mot à mot.
\end{exoresolu}

\begin{attention}[Les 5 erreurs qui coûtent des points au Bac]
\begin{enumerate}
\item \textbf{Les accents oubliés.} \esp{cooperación}, \esp{globalización}, \esp{democracia}, \esp{político}, \esp{países}, \esp{educación} : chaque accent manquant est pénalisé. Règle : les mots en \esp{-ción} portent toujours un accent sur la ó.
\item \textbf{Les faux amis.} \esp{actualmente} = « actuellement » (et non « effectivement ») ; \esp{una conferencia} = « une conférence » ; \esp{asistir a} = « assister à » (et non « aider », qui se dit \esp{ayudar}) ; \esp{la deuda} = « la dette » (et non « le doute », qui se dit \esp{la duda}).
\item \textbf{Les calques du français.} « Les droits de l'homme » se dit \esp{los derechos humanos} ; « avoir le droit de » se dit \esp{tener derecho a} (préposition \esp{a}, jamais \esp{de}) ; « selon » se dit \esp{según}, jamais \esp{segun} sans accent.
\item \textbf{Les accords et le genre.} \esp{la deuda externa} (féminin), \esp{las relaciones internacionales} (double accord), \esp{los países desarrollados}. Un adjectif non accordé dans une rédaction coûte cher.
\item \textbf{Le choix du temps et du mode.} Après \esp{no creo que}, \esp{es importante que}, le \textbf{subjonctif} est obligatoire : \esp{No creo que la ayuda sea suficiente.} Pour un souhait ou une recommandation, utilise le conditionnel : \esp{Los países ricos deberían cooperar más.} Dans le résumé d'un texte, privilégie le présent de vérité générale.
\end{enumerate}
\end{attention}

\begin{experience}[Débat citoyen : ¿Ayuda o comercio justo?]
Par groupes de quatre, préparez un débat en espagnol. Deux élèves défendent la position : \esp{Los países del Norte deben aumentar la ayuda humanitaria.} Deux autres défendent : \esp{Lo que necesita el Sur es un comercio justo, no limosna.} Chaque groupe doit utiliser au moins cinq mots du lexique de la leçon (\esp{cooperación}, \esp{deuda}, \esp{desarrollo}, \esp{derechos humanos}, \esp{globalización}) et deux connecteurs (\esp{sin embargo}, \esp{por lo tanto}). La classe vote ensuite pour la position la plus convaincante.
\end{experience}

\begin{savaistu}[Le Cameroun, la Guinée équatoriale et l'espagnol]
La Guinée équatoriale, voisine du Cameroun, est le seul pays d'Afrique subsaharienne dont l'espagnol est langue officielle : c'est un héritage de la colonisation espagnole, terminée en 1968. C'est pourquoi l'espagnol est une langue stratégique pour le Cameroun, qui entretient avec ce pays des relations commerciales et diplomatiques étroites. La Guinée équatoriale est aussi membre de l'Union africaine (\esp{la Unión Africana}), organisation continentale dont le siège est à Addis-Abeba, en Éthiopie.
\end{savaistu}

\newpage

\coursec{Exercices}

\courssub{Je vérifie mes connaissances}

\begin{exercice}[QCM : vocabulaire et compréhension]
Pour chaque question, choisis la bonne réponse (a, b ou c).

1. \esp{La ONU} signifie :
a) la Organización de Naciones Unidas \quad b) la Oficina Nacional Unida \quad c) la Organización Norte Unida

2. \esp{La deuda externa} désigne :
a) la dette intérieure d'un État \quad b) la dette d'un pays envers l'étranger \quad c) le budget de l'État

3. Une \esp{monarquía parlamentaria} est un régime où :
a) le roi gouverne seul \quad b) le roi règne mais le Parlement légifère \quad c) il n'y a pas de roi

4. Le contraire de \esp{dictadura} est :
a) \esp{república} \quad b) \esp{democracia} \quad c) \esp{monarquía}

5. \esp{La cooperación Norte-Sur} concerne :
a) les échanges entre pays riches et pays en développement \quad b) les échanges entre l'Europe et l'Asie \quad c) les guerres commerciales

6. \esp{El derecho al voto} est :
a) un devoir fiscal \quad b) un droit politique \quad c) une liberté économique

7. \esp{La globalización} signifie en français :
a) la régionalisation \quad b) la mondialisation \quad c) la colonisation

8. \esp{El desarrollo sostenible} est un développement :
a) rapide \quad b) durable \quad c) économique uniquement

9. L'\esp{Unión Africana} a son siège à :
a) Addis-Abeba \quad b) Yaoundé \quad c) Nairobi

10. Dans une démocratie, \esp{la libertad de expresión} permet :
a) de voter \quad b) d'exprimer ses opinions librement \quad c) de ne pas payer d'impôts
\end{exercice}

\begin{exercice}[Vrai ou faux : justifie ta réponse]
Dis si les affirmations suivantes sont vraies ou fausses, puis justifie en une phrase en espagnol.

1. \esp{En una dictadura, los ciudadanos eligen libremente a sus gobernantes.}

2. \esp{La libertad de prensa es un pilar de la democracia.}

3. \esp{Los derechos del ciudadano no implican ningún deber.}

4. \esp{La cooperación internacional puede ayudar a reducir la pobreza en los países del Sur.}

5. \esp{España es una república presidencialista.}
\end{exercice}

\begin{exercice}[Vocabulaire : trouve le mot exact]
Complète chaque phrase avec le mot qui convient, choisi dans la liste : \esp{deuda — desarrollo — elecciones — derechos humanos — cooperación — república}.

1. \esp{Camerún es una \ldots\ldots\ldots\ldots donde el jefe del Estado es elegido.}

2. \esp{Los ciudadanos votan en las \ldots\ldots\ldots\ldots presidenciales.}

3. \esp{La \ldots\ldots\ldots\ldots internacional reúne a países ricos y pobres para luchar contra la pobreza.}

4. \esp{Muchos países africanos sufren el peso de la \ldots\ldots\ldots\ldots externa.}

5. \esp{La Declaración Universal de los \ldots\ldots\ldots\ldots fue proclamada en 1948.}

6. \esp{El \ldots\ldots\ldots\ldots económico de un país mejora la vida de sus habitantes.}
\end{exercice}

\begin{exercice}[Conjugaison à compléter]
Conjugue les verbes entre parenthèses au temps et au mode indiqués.

1. \esp{Si los ciudadanos \ldots\ldots\ldots\ldots (votar, présent de l'indicatif), la democracia funciona.}

2. \esp{Es importante que los jóvenes \ldots\ldots\ldots\ldots (conocer, présent du subjonctif) sus derechos y deberes.}

3. \esp{Ayer, la ONU \ldots\ldots\ldots\ldots (anunciar, passé simple) un nuevo programa de ayuda.}

4. \esp{El año próximo, los países del Sur \ldots\ldots\ldots\ldots (pedir, futur simple) la cancelación de la deuda.}

5. \esp{Ojalá todos los pueblos \ldots\ldots\ldots\ldots (vivir, présent du subjonctif) en paz.}
\end{exercice}

\courssub{Je m'entraîne}

\begin{exercice}[Thème : phrases à traduire]
Traduis les phrases suivantes en espagnol. Attention aux faux amis et au choix entre \esp{ser} et \esp{estar}.

1. Les relations internationales sont très importantes pour le développement de l'Afrique.

2. Dans une démocratie, les citoyens ont des droits mais aussi des devoirs.

3. Il faut que les pays riches aident les pays pauvres. (subjonctif)

4. La Guinée équatoriale est le seul pays hispanophone d'Afrique.

5. Les jeunes doivent participer à la vie politique de leur pays.
\end{exercice}

\begin{exercice}[Transformations de phrases]
Transforme les phrases selon la consigne entre parenthèses.

1. \esp{El presidente dijo : « Respetaremos los derechos humanos. »} (mets au discours indirect : \esp{El presidente dijo que\ldots})

2. \esp{Los ciudadanos eligen a sus representantes.} (mets à la voix passive : \esp{Los representantes\ldots})

3. \esp{Si tengo tiempo, participaré en el debate.} (mets la première proposition à l'imparfait du subjonctif : \esp{Si tuviera\ldots})

4. \esp{La ONU protege los derechos de los niños.} (mets à la forme négative puis à la forme interrogative)

5. \esp{Es necesario. Los países cooperan.} (relie les deux phrases avec \esp{que} et fais les changements nécessaires)
\end{exercice}

\begin{textebox}[La vida democrática]
Vivir en una \ldots\ldots\ldots\ldots significa que los ciudadanos pueden \ldots\ldots\ldots\ldots libremente para elegir a sus gobernantes. Pero la \ldots\ldots\ldots\ldots no es total : todos debemos respetar las \ldots\ldots\ldots\ldots del país. Los derechos van siempre acompañados de \ldots\ldots\ldots\ldots : pagar los impuestos, respetar a los demás y trabajar por la \ldots\ldots\ldots\ldots social.
\end{textebox}

\begin{exercice}[Texte à trous]
Complète le texte « La vida democrática » ci-dessus avec les mots de la liste : \esp{votar — leyes — libertad — deberes — democracia — paz}.
\end{exercice}

\begin{exercice}[Appariements : mots et définitions]
Associe chaque terme (1 à 6) à sa définition (a à f).

1. \esp{la monarquía} \quad 2. \esp{la dictadura} \quad 3. \esp{la república} \quad 4. \esp{el sufragio universal} \quad 5. \esp{la libertad de prensa} \quad 6. \esp{la separación de poderes}

a) régime où une personne concentre tous les pouvoirs sans élections libres

b) droit de tous les citoyens de voter

c) régime où le chef de l'État est un roi ou une reine

d) principe qui distingue le pouvoir exécutif, législatif et judiciaire

e) droit des journalistes d'informer sans censure

f) régime où le chef de l'État est élu pour une durée limitée
\end{exercice}

\courssub{Je me prépare au Bac}

\begin{textebox}[La cooperación, una esperanza para África]
En los últimos años, las relaciones Norte-Sur han cambiado mucho. Los países ricos ya no son los únicos actores de la cooperación internacional : nuevos socios, como China o los países del Golfo, invierten en África. Sin embargo, muchos ciudadanos africanos se preguntan si esta cooperación beneficia realmente a la población. La deuda externa sigue siendo un problema grave para numerosos Estados, que deben dedicar una parte importante de su presupuesto a pagarla. Por eso, algunos economistas piden que se cancele la deuda de los países más pobres. Además, la globalización ofrece oportunidades, pero también aumenta las desigualdades entre los que tienen acceso a la educación y a la tecnología y los que no. Frente a estos desafíos, la Unión Africana defiende una idea sencilla : el desarrollo del continente debe ser obra de los propios africanos, con la solidaridad de todos.
\end{textebox}

\begin{exercice}[Compréhension de texte et questions de langue]
Lis le texte « La cooperación, una esperanza para África » ci-dessus, puis réponds aux questions.

\textbf{Questions de compréhension} (réponds en espagnol, avec des phrases complètes) :

1. ¿Qué ha cambiado en las relaciones Norte-Sur según el texto ?

2. ¿Por qué la deuda externa es un problema para los Estados africanos ?

3. ¿Qué piden algunos economistas ?

4. Según la Unión Africana, ¿quiénes deben construir el desarrollo del continente ?

5. ¿Estás de acuerdo con la última frase del texto ? Justifica tu respuesta en dos o tres frases.

\textbf{Questions de langue} :

6. Trouve dans le texte deux mots de la même famille que \esp{cooperar} et \esp{desarrollar}.

7. Réécris la phrase \esp{« algunos economistas piden que se cancele la deuda »} en commençant par \esp{La deuda\ldots}

8. Traduis en français la dernière phrase du texte.
\end{exercice}

\begin{exercice}[Thème et version]
\textbf{A. Thème.} Traduis en espagnol les phrases suivantes :

1. Les droits de l'homme doivent être respectés dans tous les pays du monde.

2. Il est important que les citoyens connaissent leurs devoirs. (subjonctif)

3. Quand j'étais enfant, je ne comprenais pas la politique.

4. Si les pays du Nord coopèrent davantage, la pauvreté diminuera.

5. On dit que la liberté de la presse est le quatrième pouvoir.

\textbf{B. Version.} Traduis en français le texte « Los derechos del ciudadano » ci-dessous (environ 80 mots).
\end{exercice}

\begin{textebox}[Los derechos del ciudadano]
Todo ciudadano tiene derechos fundamentales : el derecho a la vida, a la educación, a la salud y a expresar sus ideas libremente. Pero estos derechos solo existen de verdad si el Estado los protege y si los propios ciudadanos los defienden. En una democracia, votar no es solamente un derecho : es también un deber. El que no vota no puede quejarse después de las decisiones de los gobernantes. Por eso, los jóvenes deben interesarse por la vida pública de su país desde la escuela.
\end{textebox}

\begin{textebox}[Jóvenes y política]
Muchos jóvenes dicen que la política no les interesa. Sin embargo, las decisiones de los gobiernos afectan directamente su vida : la educación, el empleo, el medio ambiente. Un ciudadano informado es un ciudadano libre. Participar en la vida política no significa solamente votar : también es debatir, asociarse, defender causas justas y respetar las opiniones de los demás. La democracia no es un regalo : es una conquista que cada generación debe proteger.
\end{textebox}

\begin{exercice}[Rédaction : commentaire de texte]
Lis l'extrait « Jóvenes y política » ci-dessus, puis rédige en espagnol un commentaire structuré d'environ \textbf{180 mots} (introduction, deux ou trois idées développées, conclusion avec ton avis personnel). Respecte la méthode vue en classe : présente le texte, dégage le thème, analyse les arguments, donne ton point de vue.

\textbf{Consignes :} rédige un texte cohérent, utilise au moins trois connecteurs logiques (\esp{sin embargo, además, por eso, en conclusión\ldots}), un verbe au subjonctif et le vocabulaire de la leçon.
\end{exercice}



\newpage

\coursec{Corrigés des exercices}

\coursec{Leçon E13 : Lectura y comentario — Relaciones internacionales, Norte-Sur, sistemas políticos, derechos y deberes}

\courssub{Je vérifie mes connaissances}

\begin{corrige}[Exercice 1 — QCM : vocabulaire et compréhension]
1. \textbf{a)} \esp{la Organización de las Naciones Unidas}. \quad 2. \textbf{b)} la dette d'un pays envers l'étranger. \quad 3. \textbf{b)} le roi règne mais le Parlement légifère (c'est le cas de l'Espagne). \quad 4. \textbf{b)} \esp{democracia}. \quad 5. \textbf{a)} les échanges entre pays riches et pays en développement. \quad 6. \textbf{b)} un droit politique. \quad 7. \textbf{b)} la mondialisation. \quad 8. \textbf{b)} durable. \quad 9. \textbf{a)} Addis-Abeba (Éthiopie). \quad 10. \textbf{b)} d'exprimer ses opinions librement.

\textbf{Points de vigilance :} ne pas confondre \esp{república} et \esp{democracia} : une république est une forme d'État (chef de l'État élu), la démocratie est un régime fondé sur la souveraineté du peuple. Attention au faux ami : \esp{la deuda} = la dette (et non « le deuil », qui se dit \esp{el luto}).
\end{corrige}

\begin{corrige}[Exercice 2 — Vrai ou faux]
1. \textbf{Falso.} \esp{En una dictadura, una persona o un grupo concentra el poder sin elecciones libres.}

2. \textbf{Verdadero.} \esp{La libertad de prensa permite informar a los ciudadanos y controlar a los gobernantes.}

3. \textbf{Falso.} \esp{Los derechos van siempre acompañados de deberes, como respetar las leyes y pagar los impuestos.}

4. \textbf{Verdadero.} \esp{La cooperación internacional financia escuelas, hospitales y carreteras en los países del Sur.}

5. \textbf{Falso.} \esp{España es una monarquía parlamentaria: el jefe del Estado es el rey.}

\textbf{Points de vigilance :} la justification doit être une phrase complète avec sujet et verbe conjugué. On écrit \esp{Verdadero} / \esp{Falso} (et non \esp{cierto / falso} dans ce contexte d'exercice). En espagnol, pas d'espace avant les deux-points.
\end{corrige}

\begin{corrige}[Exercice 3 — Vocabulaire : trouve le mot exact]
1. \esp{república} \quad 2. \esp{elecciones} \quad 3. \esp{cooperación} \quad 4. \esp{deuda} \quad 5. \esp{Derechos Humanos} \quad 6. \esp{desarrollo}

Phrases complètes : \esp{Camerún es una república donde el jefe del Estado es elegido.} / \esp{Los ciudadanos votan en las elecciones presidenciales.} / \esp{La cooperación internacional reúne a países ricos y pobres para luchar contra la pobreza.} / \esp{Muchos países africanos sufren el peso de la deuda externa.} / \esp{La Declaración Universal de los Derechos Humanos fue proclamada en 1948.} / \esp{El desarrollo económico de un país mejora la vida de sus habitantes.}

\textbf{Point de vigilance :} on dit \esp{Camerún} sans article en espagnol standard (\esp{Camerún es un país africano}), avec un accent sur la dernière syllabe. \esp{Derechos Humanos} s'écrit avec des majuscules dans les textes officiels.
\end{corrige}

\begin{corrige}[Exercice 4 — Conjugaison à compléter]
1. \esp{votan} — présent de l'indicatif, 3\textsuperscript{e} personne du pluriel, verbe régulier en \esp{-ar}. \quad 2. \esp{conozcan} — subjonctif présent de \esp{conocer}, irrégulier à la 1\textsuperscript{re} personne (\esp{conozco}) et à toutes les personnes du subjonctif : \esp{conozca, conozcas, conozca, conozcamos, conozcáis, conozcan}. \quad 3. \esp{anunció} — passé simple (\esp{pretérito indefinido}), 3\textsuperscript{e} personne du singulier, avec accent obligatoire. \quad 4. \esp{pedirán} — futur simple : infinitif + \esp{-án} ; le changement de voyelle \esp{e $\rightarrow$ i} n'existe pas au futur. \quad 5. \esp{vivan} — subjonctif présent après \esp{ojalá}.

\textbf{Points de vigilance :} \esp{ojalá} est toujours suivi du subjonctif ; le futur se construit sur l'infinitif entier, jamais sur le radical du présent.
\end{corrige}

\courssub{Je m'entraîne}

\begin{corrige}[Exercice 5 — Thème : phrases à traduire]
1. \esp{Las relaciones internacionales son muy importantes para el desarrollo de África.} (caractéristique permanente : \esp{ser} ; « pour » de finalité : \esp{para}, jamais \esp{por}.)

2. \esp{En una democracia, los ciudadanos tienen derechos pero también deberes.}

3. \esp{Es necesario que los países ricos ayuden a los países pobres.} (subjonctif obligatoire après \esp{es necesario que} ; \esp{ayudar a} + personne.)

4. \esp{Guinea Ecuatorial es el único país hispanohablante de África.}

5. \esp{Los jóvenes deben participar en la vida política de su país.} (\esp{participar en}, et non \esp{participar a} : ne pas calquer « participer à ».)

\textbf{Points de vigilance :} « important » se dit \esp{importante} (un seul \esp{n}) ; « le développement » : \esp{el desarrollo} (double \esp{r}) ; \esp{hispanohablante} s'écrit en un seul mot.
\end{corrige}

\begin{corrige}[Exercice 6 — Transformations de phrases]
1. \esp{El presidente dijo que respetarían los derechos humanos.} (futur du discours direct $\rightarrow$ conditionnel au discours indirect : concordance des temps.)

2. \esp{Los representantes son elegidos por los ciudadanos.} (passif : \esp{ser} + participe passé accordé avec le sujet ; le complément d'agent est introduit par \esp{por}.)

3. \esp{Si tuviera tiempo, participaría en el debate.} (\esp{si} + imparfait du subjonctif $\rightarrow$ conditionnel dans la principale.)

4. Forme négative : \esp{La ONU no protege los derechos de los niños.} Forme interrogative : \esp{¿Protege la ONU los derechos de los niños?}

5. \esp{Es necesario que los países cooperen.} (subjonctif obligatoire après \esp{es necesario que}.)

\textbf{Points de vigilance :} ne jamais mettre le futur ou le conditionnel après \esp{si} ; à l'interrogatif, ne pas oublier le point d'interrogation ouvrant \esp{¿}.
\end{corrige}

\begin{corrige}[Exercice 7 — Texte à trous]
\esp{Vivir en una \textbf{democracia} significa que los ciudadanos pueden \textbf{votar} libremente para elegir a sus gobernantes. Pero la \textbf{libertad} no es total: todos debemos respetar las \textbf{leyes} del país. Los derechos van siempre acompañados de \textbf{deberes}: pagar los impuestos, respetar a los demás y trabajar por la \textbf{paz} social.}

\textbf{Points de vigilance :} après \esp{poder}, on emploie l'infinitif (\esp{votar}) ; \esp{la ley} a pour pluriel \esp{las leyes} ; \esp{la paz} prend un \esp{z} au singulier. Pas d'espace avant les deux-points en espagnol.
\end{corrige}

\begin{corrige}[Exercice 8 — Appariements : mots et définitions]
1 $\rightarrow$ c (\esp{la monarquía} : le chef de l'État est un roi ou une reine). \quad 2 $\rightarrow$ a (\esp{la dictadura} : concentration du pouvoir sans élections libres). \quad 3 $\rightarrow$ f (\esp{la república} : chef de l'État élu pour une durée limitée). \quad 4 $\rightarrow$ b (\esp{el sufragio universal} : droit de vote de tous les citoyens). \quad 5 $\rightarrow$ e (\esp{la libertad de prensa} : informer sans censure). \quad 6 $\rightarrow$ d (\esp{la separación de poderes} : exécutif, législatif, judiciaire).

\textbf{Point de vigilance :} \esp{el sufragio} désigne le droit de vote ; ne pas confondre avec le verbe \esp{sufragar} (« financer ») ni avec le français « suffrage » au sens de « voix ».
\end{corrige}

\courssub{Je me prépare au Bac}

\begin{corrige}[Exercice 9 — Compréhension de texte et questions de langue]
\textbf{Questions de compréhension.}

1. \esp{Según el texto, los países ricos ya no son los únicos actores: nuevos socios, como China o los países del Golfo, invierten en África.}

2. \esp{La deuda externa es un problema porque los Estados deben dedicar una parte importante de su presupuesto a pagarla.}

3. \esp{Algunos economistas piden que se cancele la deuda de los países más pobres.}

4. \esp{Según la Unión Africana, el desarrollo del continente debe ser obra de los propios africanos.}

5. Réponse personnelle, par exemple : \esp{Sí, estoy de acuerdo: los africanos conocen mejor los problemas de su continente. Sin embargo, la solidaridad internacional sigue siendo necesaria para financiar la educación y la salud.}

\textbf{Questions de langue.}

6. \esp{cooperar} $\rightarrow$ \esp{la cooperación} ; \esp{desarrollar} $\rightarrow$ \esp{el desarrollo}.

7. \esp{La deuda de los países más pobres debe ser cancelada, según algunos economistas.}

8. \esp{Ante estos desafíos, la Unión Africana defiende una idea sencilla: el desarrollo del continente debe ser obra de los propios africanos, con la solidaridad de todos.}

\textbf{Barème indicatif (sur 10) :} questions 1 à 4 : 1 point chacune (forme + contenu) ; question 5 : 2 points (avis + justification) ; questions 6 à 8 : 1 point chacune ; qualité globale de la langue : 1 point.

\textbf{Points de vigilance :} répondre par des phrases complètes en reprenant les mots de la question ; \esp{se cancele} est un subjonctif à valeur passive (\esp{que la deuda sea cancelada}).
\end{corrige}

\begin{corrige}[Exercice 10 — Thème et version]
\textbf{A. Thème.}

1. \esp{Los derechos humanos deben ser respetados en todos los países del mundo.} (passif : \esp{ser} + participe passé accordé.)

2. \esp{Es importante que los ciudadanos conozcan sus deberes.} (subjonctif après \esp{es importante que} ; \esp{conocer} : \esp{conozcan}.)

3. \esp{Cuando era niño, no entendía la política.} (imparfait pour les deux verbes : état et habitude dans le passé.)

4. \esp{Si los países del Norte cooperan más, la pobreza disminuirá.} (\esp{si} + présent de l'indicatif $\rightarrow$ futur ; jamais de futur après \esp{si}.)

5. \esp{Se dice que la libertad de prensa es el cuarto poder.} (\esp{se dice que} = « on dit que ».)

\textbf{B. Version.}

\esp{Todo ciudadano posee derechos fundamentales: el derecho a la vida, a la educación, a la salud y a expresar libremente sus ideas. Pero estos derechos solo existen realmente si el Estado los protege y si los propios ciudadanos los defienden. En una democracia, votar no es solo un derecho: es también un deber. Quien no vota no puede quejarse después de las decisiones de los gobernantes. Por eso los jóvenes deben interesarse por la vida pública de su país desde la escuela.}

\textbf{Barème indicatif (sur 10) :} thème : 5 points (1 point par phrase) ; version : 5 points (fidélité au sens : 3 points ; qualité du français : 2 points).

\textbf{Points de vigilance :} \esp{quejarse} = se plaindre (verbe pronominal) ; \esp{los gobernantes} = les gouvernants, et non « les gouverneurs » ; \esp{desde la escuela} = dès l'école ; \esp{interesarse por} (et non \esp{interesarse en}).
\end{corrige}

\begin{corrige}[Exercice 11 — Rédaction : commentaire de texte]
\textbf{Proposition de commentaire modèle :}

\esp{El texto « Jóvenes y política » trata de la participación de los jóvenes en la vida democrática. El autor parte de una observación: muchos jóvenes dicen que la política no les interesa. Sin embargo, defiende la idea de que participar es una necesidad para todos los ciudadanos.}

\esp{En primer lugar, el texto recuerda que las decisiones de los gobiernos afectan directamente la vida de los jóvenes: la educación, el empleo y el medio ambiente son temas políticos. Además, el autor amplía la idea de participación: no significa solamente votar, sino también debatir, asociarse y defender causas justas. Por eso, afirma que un ciudadano informado es un ciudadano libre.}

\esp{En mi opinión, el texto tiene razón: es importante que los jóvenes conozcan sus derechos y sus deberes desde la escuela. En Camerún, por ejemplo, los estudiantes pueden organizar debates en sus liceos para comprender mejor la vida pública.}

\esp{En conclusión, la democracia no es un regalo sino una conquista que cada generación debe proteger. Los jóvenes deben interesarse por la política si quieren construir un futuro más justo.} (environ 170 mots.)

\textbf{Barème indicatif (sur 20) :} respect de la structure (introduction, développement, conclusion) : 4 points ; pertinence des idées et compréhension du texte : 6 points ; correction grammaticale (accords, temps, subjonctif) : 5 points ; richesse du vocabulaire et connecteurs logiques : 3 points ; avis personnel argumenté : 2 points.

\textbf{Points de vigilance :} employer au moins trois connecteurs (\esp{sin embargo, además, por eso, en conclusión}) ; vérifier la présence d'un subjonctif (\esp{es importante que los jóvenes conozcan\ldots}) ; éviter les calques du français : on écrit \esp{interesarse por la política} et \esp{participar en}, jamais \esp{participar a}.
\end{corrige}

\newpage

\coursec{Fiche bilan}

\begin{popfigure}
\begin{center}
\begin{tikzpicture}[mindmap, grow cyclic, every node/.style={concept, font=\small}, root concept/.append style={concept color=popblue, font=\bfseries}, level 1/.append style={level distance=3.6cm, sibling angle=51}, text=white]
\node[root concept]{Relaciones internacionales}
  child[concept color=poporange]{ node {Lexique : cooperación, ONU, deuda, globalización} }
  child[concept color=popgreen]{ node {Systèmes : democracia, dictadura, monarquía, república} }
  child[concept color=poppurple]{ node {Droits et devoirs du citoyen} }
  child[concept color=poppink]{ node {Grammaire : subjonctif d'opinion et futur} }
  child[concept color=popgold]{ node {Méthode du commentaire de texte} }
  child[concept color=popmuted]{ node {Relations Nord-Sud : desarrollo, ayuda} }
  child[concept color=popdark]{ node {Production : prise de position citoyenne} };
\end{tikzpicture}
\end{center}
\legende{Carte mentale de la leçon : relations internationales, systèmes politiques, droits et devoirs.}
\end{popfigure}

\begin{bilan}
\textbf{1. Lexique clé à connaître par cœur}
\begin{center}
\begin{tabularx}{\linewidth}{|X|X|}
\hline
\rowcolor{popblueL} \textbf{Español} & \textbf{Français} \\
\hline
las relaciones internacionales & les relations internationales \\
\hline
la cooperación / el desarrollo & la coopération / le développement \\
\hline
la ONU (Organización de las Naciones Unidas) & l'ONU \\
\hline
la Unión Africana & l'Union africaine \\
\hline
la deuda externa & la dette extérieure \\
\hline
la globalización & la mondialisation \\
\hline
la ayuda al desarrollo & l'aide au développement \\
\hline
la democracia / la dictadura & la démocratie / la dictature \\
\hline
la monarquía / la república & la monarchie / la république \\
\hline
los derechos humanos & les droits de l'homme \\
\hline
los deberes del ciudadano & les devoirs du citoyen \\
\hline
la libertad de expresión & la liberté d'expression \\
\hline
el derecho al voto & le droit de vote \\
\hline
el respeto de las leyes & le respect des lois \\
\hline
\end{tabularx}
\end{center}

\textbf{2. Règles de grammaire indispensables}
\begin{center}
\begin{tabularx}{\linewidth}{|X|X|}
\hline
\rowcolor{popgreenL} \textbf{Règle} & \textbf{Exemple} \\
\hline
Subjonctif après les verbes d'opinion négative ou de doute : \esp{no creo que}, \esp{es posible que} & \esp{No creo que la deuda sea justa.} « Je ne crois pas que la dette soit juste. » \\
\hline
Subjonctif après \esp{es importante que}, \esp{es necesario que} & \esp{Es necesario que los países cooperen.} \\
\hline
Futur simple pour les prévisions et les promesses politiques & \esp{Los jóvenes votarán mañana.} \\
\hline
Conditionnel pour l'hypothèse et le conseil & \esp{Sería bueno respetar los derechos humanos.} \\
\hline
\esp{hay que} + infinitif pour le devoir général & \esp{Hay que respetar las leyes.} \\
\hline
\end{tabularx}
\end{center}

\textbf{3. Méthodes express}
\begin{itemize}
  \item \textbf{Commentaire de texte} : introduction (présenter le texte : nature, thème, problématique) ; développement (idées principales, arguments, lexique du champ politique) ; conclusion (bilan et ouverture).
  \item \textbf{Prise de position} : annoncer son opinion (\esp{en mi opinión}, \esp{pienso que}), justifier avec \esp{porque}, \esp{ya que}, nuancer avec \esp{sin embargo}, conclure avec \esp{en conclusión}.
  \item \textbf{Vérifier} : accords en genre et en nombre, subjonctif après les expressions de doute, accents écrits.
\end{itemize}
\end{bilan}

\begin{aretenir}[Checklist avant le Bac]
\begin{itemize}
  \item $\square$ Je connais le lexique des relations internationales (\esp{cooperación}, \esp{deuda}, \esp{globalización}, \esp{desarrollo}).
  \item $\square$ Je sais nommer les principales organisations : \esp{la ONU}, \esp{la Unión Africana}.
  \item $\square$ Je distingue \esp{democracia}, \esp{dictadura}, \esp{monarquía} et \esp{república}.
  \item $\square$ Je sais citer au moins trois droits du citoyen en espagnol.
  \item $\square$ Je sais citer au moins trois devoirs du citoyen en espagnol.
  \item $\square$ J'emploie le subjonctif après \esp{no creo que} et \esp{es necesario que}.
  \item $\square$ J'utilise \esp{hay que} + infinitif pour exprimer le devoir.
  \item $\square$ Je sais construire la problématique d'un commentaire de texte.
  \item $\square$ Je maîtrise les connecteurs d'opinion : \esp{en mi opinión}, \esp{sin embargo}, \esp{en conclusión}.
  \item $\square$ Je vérifie les accords et les accents dans ma production écrite.
\end{itemize}
\end{aretenir}

\findecours

\end{document}
